% Options for packages loaded elsewhere
\PassOptionsToPackage{unicode}{hyperref}
\PassOptionsToPackage{hyphens}{url}
%
\documentclass[
  10pt,
  a4paper,
]{article}
\usepackage{amsmath,amssymb}
\usepackage{iftex}
\ifPDFTeX
  \usepackage[T1]{fontenc}
  \usepackage[utf8]{inputenc}
  \usepackage{textcomp} % provide euro and other symbols
\else % if luatex or xetex
  \usepackage{unicode-math} % this also loads fontspec
  \defaultfontfeatures{Scale=MatchLowercase}
  \defaultfontfeatures[\rmfamily]{Ligatures=TeX,Scale=1}
\fi
\usepackage{lmodern}
\ifPDFTeX\else
  % xetex/luatex font selection
\fi
% Use upquote if available, for straight quotes in verbatim environments
\IfFileExists{upquote.sty}{\usepackage{upquote}}{}
\IfFileExists{microtype.sty}{% use microtype if available
  \usepackage[]{microtype}
  \UseMicrotypeSet[protrusion]{basicmath} % disable protrusion for tt fonts
}{}
\makeatletter
\@ifundefined{KOMAClassName}{% if non-KOMA class
  \IfFileExists{parskip.sty}{%
    \usepackage{parskip}
  }{% else
    \setlength{\parindent}{0pt}
    \setlength{\parskip}{6pt plus 2pt minus 1pt}}
}{% if KOMA class
  \KOMAoptions{parskip=half}}
\makeatother
\usepackage{xcolor}
\usepackage[top=22mm,bottom=24mm,left=22mm,right=22mm]{geometry}
\usepackage{color}
\usepackage{fancyvrb}
\newcommand{\VerbBar}{|}
\newcommand{\VERB}{\Verb[commandchars=\\\{\}]}
\DefineVerbatimEnvironment{Highlighting}{Verbatim}{commandchars=\\\{\}}
% Add ',fontsize=\small' for more characters per line
\newenvironment{Shaded}{}{}
\newcommand{\AlertTok}[1]{\textcolor[rgb]{1.00,0.00,0.00}{\textbf{#1}}}
\newcommand{\AnnotationTok}[1]{\textcolor[rgb]{0.38,0.63,0.69}{\textbf{\textit{#1}}}}
\newcommand{\AttributeTok}[1]{\textcolor[rgb]{0.49,0.56,0.16}{#1}}
\newcommand{\BaseNTok}[1]{\textcolor[rgb]{0.25,0.63,0.44}{#1}}
\newcommand{\BuiltInTok}[1]{\textcolor[rgb]{0.00,0.50,0.00}{#1}}
\newcommand{\CharTok}[1]{\textcolor[rgb]{0.25,0.44,0.63}{#1}}
\newcommand{\CommentTok}[1]{\textcolor[rgb]{0.38,0.63,0.69}{\textit{#1}}}
\newcommand{\CommentVarTok}[1]{\textcolor[rgb]{0.38,0.63,0.69}{\textbf{\textit{#1}}}}
\newcommand{\ConstantTok}[1]{\textcolor[rgb]{0.53,0.00,0.00}{#1}}
\newcommand{\ControlFlowTok}[1]{\textcolor[rgb]{0.00,0.44,0.13}{\textbf{#1}}}
\newcommand{\DataTypeTok}[1]{\textcolor[rgb]{0.56,0.13,0.00}{#1}}
\newcommand{\DecValTok}[1]{\textcolor[rgb]{0.25,0.63,0.44}{#1}}
\newcommand{\DocumentationTok}[1]{\textcolor[rgb]{0.73,0.13,0.13}{\textit{#1}}}
\newcommand{\ErrorTok}[1]{\textcolor[rgb]{1.00,0.00,0.00}{\textbf{#1}}}
\newcommand{\ExtensionTok}[1]{#1}
\newcommand{\FloatTok}[1]{\textcolor[rgb]{0.25,0.63,0.44}{#1}}
\newcommand{\FunctionTok}[1]{\textcolor[rgb]{0.02,0.16,0.49}{#1}}
\newcommand{\ImportTok}[1]{\textcolor[rgb]{0.00,0.50,0.00}{\textbf{#1}}}
\newcommand{\InformationTok}[1]{\textcolor[rgb]{0.38,0.63,0.69}{\textbf{\textit{#1}}}}
\newcommand{\KeywordTok}[1]{\textcolor[rgb]{0.00,0.44,0.13}{\textbf{#1}}}
\newcommand{\NormalTok}[1]{#1}
\newcommand{\OperatorTok}[1]{\textcolor[rgb]{0.40,0.40,0.40}{#1}}
\newcommand{\OtherTok}[1]{\textcolor[rgb]{0.00,0.44,0.13}{#1}}
\newcommand{\PreprocessorTok}[1]{\textcolor[rgb]{0.74,0.48,0.00}{#1}}
\newcommand{\RegionMarkerTok}[1]{#1}
\newcommand{\SpecialCharTok}[1]{\textcolor[rgb]{0.25,0.44,0.63}{#1}}
\newcommand{\SpecialStringTok}[1]{\textcolor[rgb]{0.73,0.40,0.53}{#1}}
\newcommand{\StringTok}[1]{\textcolor[rgb]{0.25,0.44,0.63}{#1}}
\newcommand{\VariableTok}[1]{\textcolor[rgb]{0.10,0.09,0.49}{#1}}
\newcommand{\VerbatimStringTok}[1]{\textcolor[rgb]{0.25,0.44,0.63}{#1}}
\newcommand{\WarningTok}[1]{\textcolor[rgb]{0.38,0.63,0.69}{\textbf{\textit{#1}}}}
\usepackage{longtable,booktabs,array}
\usepackage{calc} % for calculating minipage widths
% Correct order of tables after \paragraph or \subparagraph
\usepackage{etoolbox}
\makeatletter
\patchcmd\longtable{\par}{\if@noskipsec\mbox{}\fi\par}{}{}
\makeatother
% Allow footnotes in longtable head/foot
\IfFileExists{footnotehyper.sty}{\usepackage{footnotehyper}}{\usepackage{footnote}}
\makesavenoteenv{longtable}
\setlength{\emergencystretch}{3em} % prevent overfull lines
\providecommand{\tightlist}{%
  \setlength{\itemsep}{0pt}\setlength{\parskip}{0pt}}
\setcounter{secnumdepth}{-\maxdimen} % remove section numbering
\ifLuaTeX
\usepackage[bidi=basic]{babel}
\else
\usepackage[bidi=default]{babel}
\fi
\babelprovide[main,import]{american}
% get rid of language-specific shorthands (see #6817):
\let\LanguageShortHands\languageshorthands
\def\languageshorthands#1{}
% Additional layout and mathematical packages for the EIT theory documents.
\usepackage{mathtools}
\usepackage{enumitem}
\setlist{nosep}
\allowdisplaybreaks[3]
\setlength{\emergencystretch}{3em}
\IfFileExists{xurl.sty}{\usepackage{xurl}}{}
\urlstyle{same}
\ifLuaTeX
  \usepackage{selnolig}  % disable illegal ligatures
\fi
\usepackage{bookmark}
\IfFileExists{xurl.sty}{\usepackage{xurl}}{} % add URL line breaks if available
\urlstyle{same}
\hypersetup{
  pdftitle={Generalized EIT No-Go/Go Theory Applicable to Arbitrary Materials},
  pdflang={en-US},
  hidelinks,
  pdfcreator={LaTeX via pandoc}}

\title{Generalized EIT No-Go/Go Theory Applicable to Arbitrary
Materials}
\author{}
\date{}

\begin{document}
\maketitle

\textbf{Integrated version:} v6.2 --- sector-resolved finite-dimensional
Markovian weak-probe formulation\\
\textbf{Updated:} 2026-07-13\\
\textbf{Objective:} To determine whether a specified EIT channel in a
finite-dimensional Markovian quantum system---including multiple excited
levels, internal-state mixing, phonons, dissipation, external fields,
polarization, and static inhomogeneity---belongs to the exact
sector-resolved no-go, asymptotic-suppression, kernel-go,
practical-no-go, or observable-go class in the stationary rotating-frame
weak-probe regime.

\begin{center}\rule{0.5\linewidth}{0.5pt}\end{center}

\subsection{0. Executive summary}\label{executive-summary}

This theory does not classify EIT from the material name alone. The
object to be classified is, at minimum,

\[
\boxed{
\mathfrak C=
\{\text{material},\ \mathcal H_g,\ \mathcal H_e,\ \text{lower-state pair},\
\text{probe/control polarization},\ \text{external fields},\ T,\ \text{observable},\ \mathbb S\}
}
\]

a \textbf{material- and channel-resolved configuration} specified by the
above tuple. Even within the same material, a spin-\(\Lambda\) channel,
an orbital-\(\Lambda\) channel, different polarizations, and different
magnetic-field or strain conditions are distinct classification targets.

The theory rigorously separates the following four concepts.

\begin{enumerate}
\def\labelenumi{\arabic{enumi}.}
\tightlist
\item
  \textbf{Optical dark subspace}\\
  Instantaneous optical decoupling is determined by the nullspace of the
  coupling map \(\Omega\).
\item
  \textbf{Stationary Lindblad dark state}\\
  For a pure dark state to be stationary, additional conditions on the
  lower-manifold dynamics and on every Lindblad jump operator are
  required.
\item
  \textbf{Coherent EIT pathway}\\
  In the reduced model, the excited-manifold resolvent kernel \[
  \mathbf K(z)=C^\dagger A(z)^{-1}C
  \] has off-diagonal elements that determine the feedback between the
  probe optical sector and the long-lived coherence sector.
\item
  \textbf{Observable EIT}\\
  This is determined by the response including the full Liouvillian,
  ground-state coherence, populations, control power, optical depth,
  inhomogeneous broadening, propagation, and the detection scheme.
\end{enumerate}

In this version, a no-go condition in the full Liouvillian is not
defined as a zero of the total susceptibility. We fix a physically
specified long-lived sector \(\mathbb S\) and its complementary
direct-response sector \(\mathbb X\). Let
\(\chi_{\rm cut}^{(\mathbb S)}\) denote the counterfactual response
obtained by cutting only the cross-feedback with \(\mathbb S\), and let
\(\chi_{\rm full}\) denote the complete response. We then define

\[
\boxed{
\delta\chi_{\mathbb S}
\equiv
\chi_{\rm full}-\chi_{\rm cut}^{(\mathbb S)}
}
\]

as the \textbf{\(\mathbb S\)-sector-mediated response modification}.
When \(\mathbb S\) is a long-lived off-diagonal coherence sector and the
cross-coupling is generated by the control Hamiltonian, this quantity
may be interpreted as the coherent EIT correction. A generic Liouvillian
block can also contain dissipative coupling, so in that case it should
not be called merely a coherent correction.

A structural exact EIT no-go condition is that, for the specified sector
partition and parameter domain,

\[
\boxed{
\delta\chi_{\mathbb S}(\theta)\equiv0
}
\]

holds identically. This is not an unconditional scalar label intrinsic
to the system; it is a sector-resolved statement about the sector
\(\mathbb S\) included in the configuration tuple. By contrast,

\[
\chi_{\rm full}=0
\]

may represent complete transparency, i.e., ideal EIT, and therefore is
not a no-go condition.

For a closed two-lower-state weak-probe block, taking
\(\mathbb S=\operatorname{span}\{\rho_{g_2g_1}^{(1)}\}\), one obtains at
a regular point

\[
\delta\Xi_{\mathbb S}
=-\frac{\beta K_{12}K_{21}}
{\gamma_g+\beta S_2}.
\]

In this scalar block, the cross-coupling originates from the control
Hamiltonian, so \(\delta\Xi_{\mathbb S}\) coincides with the coherent
EIT correction.

Therefore, for nonzero control and a finite denominator,

\[
\boxed{K_{12}K_{21}=0}
\]

is the necessary and sufficient zero condition for the targeted
coherence feedback.

Furthermore,

\begin{itemize}
\tightlist
\item
  Under strong full-rank damping, the first nonzero resolvent moment
  determines the suppression exponent.
\item
  An exact zero results if all moments vanish, or if a
  generator-preserving reducing symmetry separates the target transfer.
\item
  If the damping operator is singular, an \(O(1)\) channel protected in
  \(\ker D\) may remain.
\item
  If the ground-state index is mixed, the full Liouville-space response
  operator must be used instead of a Hilbert-space kernel.
\item
  Exact \(\mathbb S\)-sector no-go, asymptotic suppression, practical
  visibility, and observable EIT must be reported separately.
\end{itemize}

Accordingly, the present theory does not mean that ``all arbitrary
materials have already been classified automatically.'' Its scope is an
\textbf{EIT channel that possesses a Markovian GKSL generator on a
finite-dimensional Hilbert/Liouville space and admits a weak-probe
linear response in a stationary rotating frame}. Infinite-dimensional
continua, genuinely non-Markovian memory, strong-probe nonlinear
response, and propagation-dominated collective EIT lie outside the scope
of the main theorem.

\section{1. Scope and Unit of the
Claim}\label{scope-and-unit-of-the-claim}

\subsection{1.1 Classifying channels rather than
materials}\label{classifying-channels-rather-than-materials}

The no-go/go classification is defined for each configuration tuple
below.

\[
\mathfrak C=
(M,G,E,a,b,\epsilon_p,\epsilon_c,\mathbf F,T,\mathcal O,\mathbb S),
\]

Here,

\begin{itemize}
\tightlist
\item
  \(M\): material or emitter
\item
  \(G\subset\mathcal H_g\): lower manifold used in the channel
\item
  \(E\subset\mathcal H_e\): operative excited manifold
\item
  \(a,b\): lower-state pair
\item
  \(\epsilon_p,\epsilon_c\): probe/control polarization
\item
  \(\mathbf F\): magnetic field, electric field, strain, cavity, etc.
\item
  \(T\): temperature or bath condition
\item
  \(\mathcal O\): observable such as absorption, transmission, or
  fluorescence
\item
  \(\mathbb S\): targeted long-lived coherence sector for the no-go/go
  classification
\end{itemize}

Thus,

the following are distinct classification targets.

\begin{itemize}
\tightlist
\item
  zero-field ZPL spin-\(\Lambda\) and transverse-field spin-\(\Lambda\)
  channels in the same NV center
\item
  spin-\(\Lambda\) and orbital-\(\Lambda\) channels in the same group-IV
  center
\item
  parallel and cross polarizations for the same lower-state pair
\item
  a resolved manifold and a thermally merged manifold in the same
  material
\end{itemize}

\subsection{1.2 Meaning of ``applicable to arbitrary
materials''}\label{meaning-of-applicable-to-arbitrary-materials}

The claim that the theory is material independent means the following.

\begin{quote}
Given the material-specific Hamiltonian, dipole matrix elements, jump
operators, bath rates, and external-field conditions, the same
mathematical classifier can evaluate the exact zero, asymptotic order,
coherent-pathway permission, and observable visibility of an EIT
channel.
\end{quote}

It does not mean the following.

\begin{itemize}
\tightlist
\item
  A no-go/go result can be returned from the material name alone.
\item
  \(K\neq0\) guarantees observable EIT.
\item
  Orbital degeneracy necessarily implies a go result.
\item
  The response automatically becomes scalar when the linewidth exceeds
  the fine-structure splitting.
\item
  The finite-dimensional moment theorem applies without modification to
  every non-Markovian continuum.
\end{itemize}

\subsection{1.3 Assumptions of the main
theorems}\label{assumptions-of-the-main-theorems}

The exact theorems and solver validation in this framework are
restricted to the following assumptions.

\begin{enumerate}
\def\labelenumi{\arabic{enumi}.}
\tightlist
\item
  \textbf{Finite dimensionality}\\
  The Hilbert space, or the adopted Liouville response space, is finite
  dimensional. When a large vibronic manifold is truncated, convergence
  of both the observable and the kernel must be checked separately.
\item
  \textbf{Markovian GKSL dynamics}\\
  The zeroth-order dynamics, including the control field, external
  fields, and dissipation, is represented in a stationary rotating frame
  by a time-independent GKSL generator \(\mathcal L_c\).
\item
  \textbf{Weak-probe linear response}\\
  The response is retained to \(O(\epsilon)\) in the probe amplitude
  \(\epsilon\). The control may be treated nonperturbatively, but
  probe-induced saturation and population redistribution are not
  included in the main theorem.
\item
  \textbf{Stationary response problem}\\
  The system has a unique steady state and a finite gap, or a
  group/Drazin inverse can be specified uniquely on a compatible
  trace-zero subspace.
\item
  \textbf{Physically fixed sector partition}\\
  Before a no-go statement is made, the long-lived sector \(\mathbb S\),
  direct-response sector \(\mathbb X\), source, and detector are
  physically specified.
\item
  \textbf{Analytic-domain claim}\\
  If an identically vanishing transfer is claimed over an entire
  parameter domain, the matrix family and source/detector must be real
  analytic, or holomorphic after an appropriate complexification.
\end{enumerate}

The following cases are outside the scope of the main theorem.

\begin{itemize}
\tightlist
\item
  Systems with a genuinely non-Markovian memory kernel for which a
  finite-dimensional Markovian embedding has not converged
\item
  arbitrary strong-probe、saturation、bistability
\item
  Systems in which propagation, cavity self-consistency, or many-body
  interactions cannot be separated from the local response
\item
  Time-periodic drives that have not been reduced to a
  finite-dimensional Floquet--Liouville generator
\end{itemize}

These cases can be treated by separate extensions, but they are not
included in the present claim of applicability to arbitrary materials.

\begin{center}\rule{0.5\linewidth}{0.5pt}\end{center}

\section{2. General Open-System
Formulation}\label{general-open-system-formulation}

\subsection{2.1 Hilbert-space
decomposition}\label{hilbert-space-decomposition}

Decompose the total Hilbert space as

\[
\mathcal H=\mathcal H_e\oplus\mathcal H_g\oplus\mathcal H_r
\]

where \(\mathcal H_e\) is the selected excited manifold,
\(\mathcal H_g\) is the lower manifold, and \(\mathcal H_r\) is an
auxiliary space containing singlets, charge states, and other loss
channels.

Define the optical-coherence block of the density operator by

\[
\sigma=P_e\rho P_g\in\mathcal B(\mathcal H_g,\mathcal H_e)
\]

.

\subsection{2.2 Action of excited-confined jumps on optical
coherence}\label{action-of-excited-confined-jumps-on-optical-coherence}

When a jump operator is confined to the excited manifold,

\[
L_\mu=P_eL_\mu P_e,
\]

the Lindblad dissipator acts on the optical-coherence block as

\[
P_e\mathcal D[L_\mu](\rho)P_g
=-\frac12L_\mu^\dagger L_\mu\,\sigma
\]

. The recycling term \(L_\mu\rho L_\mu^\dagger\) vanishes in the
rectangular \(e-g\) block.

Thus, even a jump that generates exchange among excited-state
populations acts in general as left-loss damping on an optical
coherence. This is why the present theory does not use fixed-point
projection as a general rule and instead centers the analysis on
full-rank damping and inverse-rate expansions.

\subsection{2.3 Liouville-space
generator}\label{liouville-space-generator}

The undriven optical-coherence equation can generally be written as

\[
\dot\sigma
=-i(H_e\sigma-\sigma H_g)
-\frac12R_e\sigma
+\mathcal J_{eg}[\sigma]
+\text{sources},
\]

\[
R_e=\sum_\mu L_\mu^\dagger L_\mu,
\]

. Upon vectorization,

\[
\operatorname{vec}(AXB)=(B^{\mathsf T}\otimes A)\operatorname{vec}X
\]

and hence

\[
\mathbb A_{eg}
=I_g\otimes\left(\frac12R_e+iH_e\right)
-iH_g^{\mathsf T}\otimes I_e
+\mathbb A_{\rm corr}
\]

is obtained. Here \(\mathbb A_{\rm corr}\) includes ground-state mixing,
correlated jumps, spectral-diffusion models, and related corrections.

\subsection{2.4 Conditions for a Hilbert-space
reduction}\label{conditions-for-a-hilbert-space-reduction}

If the lower state \(|g_j\rangle\) is an eigenstate of \(H_g\) and the
ground-state index \(j\) is not mixed in the linear-response
optical-coherence block,

\[
\sigma=|x_j\rangle\langle g_j|
\]

the corresponding subspace is invariant, and the problem reduces exactly
to the excited-manifold Hilbert-space resolvent

\[
A_j(z)|x_j\rangle=|s_j\rangle
\]

.

The full Liouville-space operator must be retained in the following
cases.

\begin{itemize}
\tightlist
\item
  Ground-state relaxation mixes the right index.
\item
  Correlated jumps act on both the excited and ground manifolds.
\item
  The control field strongly couples several ground-state coherences.
\item
  Population redistribution cannot be neglected.
\item
  A drive other than the probe is treated nonperturbatively.
\end{itemize}

\begin{center}\rule{0.5\linewidth}{0.5pt}\end{center}

\section{3. Optical coupling、optical dark subspace、stationary dark
state}\label{optical-couplingoptical-dark-subspacestationary-dark-state}

\subsection{3.1 Coupling map}\label{coupling-map}

Define the lower-to-excited optical coupling in the rotating frame by

\[
V_+=P_eH_{\rm drive}P_g=\frac12\Omega,
\qquad
\Omega:\mathcal H_g\rightarrow\mathcal H_e
\]

. For a lower basis \(\{|g_a\rangle\}_{a=1}^{N_g}\),

\[
\Omega=(\Omega_1d_1,\ldots,\Omega_{N_g}d_{N_g}),
\]

\[
(d_a)_j
=\langle e_j|\epsilon_a\cdot\hat{\mathbf d}|g_a\rangle.
\]

The complex phases of the dipole vectors are retained. Under the basis
transformation \(U_e\) on \(\mathcal H_e\) and \(U_g\) on
\(\mathcal H_g\),

\[
\Omega\mapsto U_e\Omega U_g^\dagger
\]

so the rank and nullity are basis invariant.

\subsection{3.2 Theorem 1A --- Optical dark-subspace
rank}\label{theorem-1a-optical-dark-subspace-rank}

\textbf{Theorem.} A lower-state vector \(|v\rangle\in\mathcal H_g\) is
instantaneously decoupled from optical excitation, namely,

\[
V_+|v\rangle=0
\]

if and only if

\[
\Omega|v\rangle=0
\]

. Therefore, the optical dark subspace

\[
\mathcal D_{\rm opt}=\ker\Omega
\]

has dimension

\[
\boxed{
\dim\mathcal D_{\rm opt}
=N_g-\operatorname{rank}\Omega
}
\]

.

\textbf{Proof.} Since \(V_+=\Omega/2\), the first equivalence follows
directly from the definition. The dimension formula follows from the
rank--nullity theorem for finite-dimensional linear maps. \(\square\)

For a two-lower-state system with both field amplitudes nonzero, a
nontrivial optical dark vector exists if and only if

\[
\operatorname{rank}(d_1,d_2)<2.
\]

In particular, if both dipole vectors are nonzero,

\[
\operatorname{rank}(d_1,d_2)=1
\]

.

\subsection{3.3 Distinction between an optical dark vector and a
Hamiltonian dark
eigenstate}\label{distinction-between-an-optical-dark-vector-and-a-hamiltonian-dark-eigenstate}

Even when \(|D\rangle\in\ker\Omega\), it is not a stationary dark state
if the lower-manifold Hamiltonian rotates the state from the dark
subspace into the bright subspace.

Write the rotating-frame Hamiltonian as

\[
H=
\begin{pmatrix}
H_e & \Omega/2\\
\Omega^\dagger/2 & H_g^{(\rm rot)}
\end{pmatrix}
\]

. A pure state \(|D\rangle\) fully supported in the lower manifold is a
Hamiltonian eigenstate if and only if

\[
\boxed{
\Omega|D\rangle=0,
\qquad
H_g^{(\rm rot)}|D\rangle=E_D|D\rangle
}
\]

.

More generally, the entire dark subspace is invariant under the
Hamiltonian if and only if, with \(P_D\) denoting the orthogonal
projector onto \(\ker\Omega\),

\[
\boxed{
(I-P_D)H_g^{(\rm rot)}P_D=0
}
\]

. In the usual two-lower-state \(\Lambda\) system, this condition
incorporates two-photon resonance and lower-state energy matching.

\subsection{3.4 Theorem 1B --- Full Lindblad stationary pure dark
state}\label{theorem-1b-full-lindblad-stationary-pure-dark-state}

Let the full master equation be

\[
\dot\rho=\mathcal L(\rho)
=-i[H,\rho]
+\sum_\mu\left(
L_\mu\rho L_\mu^\dagger
-\frac12\{L_\mu^\dagger L_\mu,\rho\}
\right)
\]

.

\textbf{Theorem.} For a normalized vector \(|D\rangle\), the pure state
\(\rho_D=|D\rangle\langle D|\) is stationary, i.e.,
\(\mathcal L(\rho_D)=0\), if and only if there exist complex numbers
\(\lambda_\mu\) and a real number \(h\) such that

\[
\boxed{
L_\mu|D\rangle=\lambda_\mu|D\rangle
\quad(\forall\mu)
}
\]

and

\[
\boxed{
\left[
H+\frac{i}{2}\sum_\mu
\left(\lambda_\mu^*L_\mu-
\lambda_\mu L_\mu^\dagger\right)
\right]|D\rangle
=h|D\rangle
}
\]

hold.

\textbf{Proof sketch.} For a stationary pure state, the time derivative
of the purity vanishes. For the Lindblad equation,

\[
\frac{d}{dt}\operatorname{Tr}(\rho_D^2)
=-2\sum_\mu
\left(
\langle L_\mu^\dagger L_\mu\rangle_D
-|\langle L_\mu\rangle_D|^2
\right)
\]

and each term is a nonnegative variance. Therefore, for every jump
operator, \(L_\mu|D\rangle=\lambda_\mu|D\rangle\) is necessary.
Substituting this condition into the master equation and projecting from
the left with \(Q_D=I-|D\rangle\langle D|\), one obtains

\[
Q_D\left[
H+\frac{i}{2}\sum_\mu
(\lambda_\mu^*L_\mu-
\lambda_\mu L_\mu^\dagger)
\right]|D\rangle=0
\]

. The operator in square brackets is Hermitian, so its action is
parallel to \(|D\rangle\) and has a real eigenvalue \(h\). Conversely,
substituting the two conditions into the master equation eliminates all
off-diagonal components; trace preservation and Hermiticity then give
\(\mathcal L(\rho_D)=0\). \(\square\)

\textbf{Application.} If the EIT dark state is required to lie in the
lower manifold, impose additionally

\[
P_e|D\rangle=0,
\qquad
\Omega|D\rangle=0
\]

. When a spontaneous-emission jump starts in an excited state, one
normally has \(L_\mu|D\rangle=0\). By contrast, if ground-state
dephasing, ground-state relaxation, microwave noise, or another process
violates the common-eigenvector condition for \(|D\rangle\), an optical
dark vector may exist while no stationary pure dark state exists.

\subsection{3.5 Fixed terminology for
classification}\label{fixed-terminology-for-classification}

This theory distinguishes the following notions.

\begin{enumerate}
\def\labelenumi{\arabic{enumi}.}
\tightlist
\item
  \textbf{optical dark capable}: \(\dim\ker\Omega>0\)
\item
  \textbf{Hamiltonian dark eigenstate}: \(\Omega|D\rangle=0\) and
  \(H_g^{(\rm rot)}|D\rangle=E_D|D\rangle\)
\item
  \textbf{stationary Lindblad dark state}: satisfies Theorem 1B
\item
  \textbf{observable EIT}: a narrow coherent transparency feature
  appears in the full linear response
\end{enumerate}

Accordingly, neither stationary EIT nor observable EIT is inferred from
the rank theorem alone.

\begin{center}\rule{0.5\linewidth}{0.5pt}\end{center}

\section{4. Basis-invariant coherent-response
kernel}\label{basis-invariant-coherent-response-kernel}

\subsection{4.1 Reduced Hilbert-space
kernel}\label{reduced-hilbert-space-kernel}

Let the optical-coherence generator be \(A(z)\), and define the Green
operator by

\[
G(z)=A(z)^{-1}
\]

. For the multichannel coupling matrix

\[
C=(d_1,\ldots,d_{N_g})
\]

, define

\[
\boxed{
\mathbf K(z)=C^\dagger G(z)C
}
\]

as the coherent-response matrix kernel.

Its components are

\[
K_{ab}(z)=d_a^\dagger G(z)d_b.
\]

The elements \(K_{ab}\) are invariant under a unitary change of basis in
the excited manifold,

\[
A\mapsto UAU^\dagger,
\qquad
d_a\mapsto Ud_a
\]

.

\subsection{4.2 Physical meaning}\label{physical-meaning}

\begin{itemize}
\tightlist
\item
  diagonal element \[
  S_a=K_{aa}=d_a^\dagger Gd_a
  \] represents the single-leg optical response or self-energy.
\item
  off-diagonal element \[
  K_{ab}=d_a^\dagger Gd_b
  \] represents the coherent Raman/interference vertex mediated by the
  excited manifold.
\end{itemize}

When adiabatic elimination is valid, the effective coherent Hamiltonian
in the lower manifold contains, schematically,

\[
\langle g_a|H_{\rm eff}|g_b\rangle
\propto-\Omega_a^*\Omega_bK_{ab}
\]

. The sign and numerical prefactor depend on the rotating-frame
convention.

\subsection{4.3 Full Liouville-space
kernel}\label{full-liouville-space-kernel}

When a Hilbert-space reduction is not possible, introduce a source
superoperator \(\mathbb B_b\) and a detection functional
\(\mathbb D_a\), and define

\[
\boxed{
\mathbb K_{ab}(z)
=\langle\!\langle \mathbb D_a|
\mathbb A(z)^{-1}
|\mathbb B_b\rangle\!\rangle
}
\]

. This transfer function is a generic source-to-detector kernel for an
arbitrary finite-dimensional Liouvillian. If it is chosen to represent
the total probe susceptibility, however, its zero may correspond to
complete transparency. The EIT-specific no-go object is the
sector-resolved correction \(\delta\chi_{\mathbb S}\) defined in Section
10.

Moment, Krylov, sector-symmetry, and adjugate tests are applied to the
factors of the targeted cross-transfer, or to an augmented
finite-dimensional transfer realization representing the full/cut
response difference. A zero of the total susceptibility is not called an
EIT no-go directly.

\subsection{4.4 A kernel zero is not equivalent to a dark
state}\label{a-kernel-zero-is-not-equivalent-to-a-dark-state}

In general,

\[
K_{12}=0
\quad\not\Leftrightarrow\quad
\ker\Omega\neq\{0\}.
\]

\begin{itemize}
\tightlist
\item
  Even when the coupling vectors are linearly independent, cancellation
  among resolvent pathways can yield \(K_{12}=0\).
\item
  Even when an exact dark state exists, dissipation or population
  dynamics may render the narrow EIT feature unobservable.
\end{itemize}

\begin{center}\rule{0.5\linewidth}{0.5pt}\end{center}

\section{5. Fast damping、rate mapping、moment
hierarchy}\label{fast-dampingrate-mappingmoment-hierarchy}

\subsection{5.1 Standard decomposition}\label{standard-decomposition}

Decompose the generator on the selected finite-dimensional response
space as

\[
A_\Gamma(z)=\Gamma D+A_0(z)
\]

.

\begin{itemize}
\tightlist
\item
  \(\Gamma>0\): selected fast-rate scale
\item
  \(D\): dimensionless fast-generator shape
\item
  \(A_0(z)\): detuning, fine structure, slow loss, fields, strain, and
  related terms
\end{itemize}

\textbf{Important:} \(\Gamma\) is not in general identical to a
population-relaxation rate quoted in the literature. The term
\(\Gamma D\) must be defined from the matrix that actually appears in
the optical-coherence equation.

\subsection{5.2 Microscopic mapping from population hopping to
optical-coherence
damping}\label{microscopic-mapping-from-population-hopping-to-optical-coherence-damping}

Define incoherent hopping between the excited branches \(|X\rangle\) and
\(|Y\rangle\) by

\[
L_{Y\leftarrow X}
=\sqrt{k_{Y\leftarrow X}}|Y\rangle\langle X|,
\qquad
L_{X\leftarrow Y}
=\sqrt{k_{X\leftarrow Y}}|X\rangle\langle Y|
\]

.

The population equations are

\[
\dot p_X
=-k_{Y\leftarrow X}p_X
+k_{X\leftarrow Y}p_Y,
\]

\[
\dot p_Y
=k_{Y\leftarrow X}p_X
-k_{X\leftarrow Y}p_Y.
\]

The population mode measured relative to equilibrium relaxes at the rate

\[
\Gamma_{\rm pop}
=k_{Y\leftarrow X}+k_{X\leftarrow Y}
\]

.

By contrast, for the optical coherences associated with an arbitrary
lower state \(|g\rangle\),

\[
\sigma_{Xg}=\langle X|\rho|g\rangle,
\qquad
\sigma_{Yg}=\langle Y|\rho|g\rangle
\]

the recycling term does not contribute, and

\[
\boxed{
\dot\sigma_{Xg}\big|_{\rm hop}
=-\frac{k_{Y\leftarrow X}}{2}\sigma_{Xg}
}
\]

\[
\boxed{
\dot\sigma_{Yg}\big|_{\rm hop}
=-\frac{k_{X\leftarrow Y}}{2}\sigma_{Yg}
}
\]

.

\textbf{Proof.} For example, for \(L_{Y\leftarrow X}\),

\[
P_eL\rho L^\dagger P_g=0,
\qquad
L^\dagger L=k_{Y\leftarrow X}|X\rangle\langle X|,
\]

so only the left action of the anticommutator remains in the \(X-g\)
block, with coefficient \(-k_{Y\leftarrow X}/2\). \(\square\)

\subsubsection{Symmetric two-branch
convention}\label{symmetric-two-branch-convention}

\[
k_{Y\leftarrow X}=k_{X\leftarrow Y}=k
\]

and suppose the literature or model defines

\[
\dot{(p_X-p_Y)}=-\Gamma_{XY}(p_X-p_Y)
\]

. Then

\[
\Gamma_{XY}=2k
\]

. Consequently, the additional damping of each optical coherence is

\[
\boxed{
\gamma_{\rm oc}^{(XY)}
=\frac{k}{2}
=\frac{\Gamma_{XY}}{4}
}
\]

.

Thus, if \(D=I\) is chosen on two equivalent branches, the scale used in
the moment expansion must be

\[
\boxed{
\Gamma=\gamma_{\rm oc}^{(XY)}=\Gamma_{XY}/4
}
\]

. Alternatively, if one sets \(\Gamma=\Gamma_{XY}\), then one must use
\(D=I/4\). These two conventions must not be mixed.

\subsubsection{Asymmetric or detailed-balance
hopping}\label{asymmetric-or-detailed-balance-hopping}

If \(k_{Y\leftarrow X}\neq k_{X\leftarrow Y}\),

\[
A_{\rm hop}
=\frac12
\begin{pmatrix}
k_{Y\leftarrow X}&0\\
0&k_{X\leftarrow Y}
\end{pmatrix}
\]

and in general this is not \(\Gamma_{\rm pop}/4\) times the identity.
Material-specific calculations must input the upward and downward rates
separately rather than replacing them by an inaccurate scalar rate.

More generally, define the total incoherent escape rate from an excited
state \(|e_j\rangle\) by

\[
r_j^{\rm out}=\sum_{m\neq j}k_{m\leftarrow j}
+\sum_{r\notin E}k_{r\leftarrow j}
\]

. The corresponding optical coherence then acquires the damping
contribution

\[
\frac12r_j^{\rm out}
\]

. Pure dephasing must be derived separately from the adopted
jump-operator normalization.

\subsection{5.3 Algebraic resolvent
reduction}\label{algebraic-resolvent-reduction}

For the exact algebraic result in this section, assume initially only
that \(D\) is invertible. With the definitions

\[
X(z)=-D^{-1}A_0(z),
\qquad
\nu=D^{-1}c
\]

,

\[
A_\Gamma
=D(\Gamma I-X)
\]

so, at every point where \(A_\Gamma\) is invertible,

\[
\boxed{
K_\Gamma(z)
=p^\dagger A_\Gamma(z)^{-1}c
=p^\dagger(\Gamma I-X(z))^{-1}\nu
}
\]

.

\subsection{5.4 Controlled large-rate
expansion}\label{controlled-large-rate-expansion}

Under the condition

\[
q(z)=\frac{\|X(z)\|}{\Gamma}
=\frac{\|D^{-1}A_0(z)\|}{\Gamma}<1
\]

,

\[
(\Gamma I-X)^{-1}
=\frac1\Gamma
\sum_{n=0}^{\infty}\frac{X^n}{\Gamma^n}
\]

. Therefore,

\[
\boxed{
K_\Gamma(z)
=\sum_{n=0}^{\infty}
\frac{M_n(z)}{\Gamma^{n+1}}
}
\]

\[
\boxed{
M_n(z)=p^\dagger X(z)^n\nu,
\qquad
\nu=D^{-1}c
}
\]

. The former notation mixed \(u\) and \(\nu\); this version consistently
uses \(\nu\).

The remainder after truncation at order \(m\),

\[
R_m
=K_\Gamma-
\sum_{n=0}^{m}\frac{M_n}{\Gamma^{n+1}}
\]

, satisfies

\[
\boxed{
|R_m|
\le
\frac{\|p\|\,\|D^{-1}\|\,\|c\|}{\Gamma}
\frac{q^{m+1}}{1-q}
}
\]

.

\subsection{5.5 Theorem 3 --- First nonzero
moment}\label{theorem-3-first-nonzero-moment}

\[
M_0=\cdots=M_{m-1}=0,
\qquad
M_m\neq0
\]

then, as \(\Gamma\to\infty\),

\[
\boxed{
K_\Gamma
=\frac{M_m}{\Gamma^{m+1}}
+O(\Gamma^{-m-2})
}
\]

. This follows directly from the Neumann expansion.

\subsection{5.6 Theorem 4 --- Exact finite-dimensional transfer zero:
complete
equivalence}\label{theorem-4-exact-finite-dimensional-transfer-zero-complete-equivalence}

Let the response-space dimension be \(N\), and assume that \(D\) is
invertible at fixed \(z\). The following statements are equivalent.

\begin{enumerate}
\def\labelenumi{\arabic{enumi}.}
\tightlist
\item
  \(M_n=p^\dagger X^n\nu=0\) for \(n=0,\ldots,N-1\)
\item
  \(p\) is orthogonal to the finite Krylov space \[
  \mathcal K_N(X,\nu)
  =\operatorname{span}\{\nu,X\nu,\ldots,X^{N-1}\nu\}
  \] .
\item
  For every polynomial \(q\), \[
  p^\dagger q(X)\nu=0
  \]
\item
  For every \(\Gamma\notin\operatorname{spec}X\), \[
  \boxed{K_\Gamma=0}
  \]
\end{enumerate}

\textbf{Proof.} Statements 1 and 2 are equivalent by definition. By the
Cayley--Hamilton theorem, \(X^N\) is a linear combination of
\(I,X,\ldots,X^{N-1}\). Hence statement 1 implies inductively that every
\(M_n\) vanishes, which yields statement 3.

Next, for \(\Gamma\notin\operatorname{spec}X\),

\[
(\Gamma I-X)^{-1}
=\frac{\operatorname{adj}(\Gamma I-X)}
{\det(\Gamma I-X)}.
\]

The adjugate \(\operatorname{adj}(\Gamma I-X)\) can be expressed as a
polynomial in \(X\) of degree at most \(N-1\). Statement 3 therefore
implies

\[
p^\dagger(\Gamma I-X)^{-1}\nu=0
\]

which gives statement 4.

Conversely, if statement 4 holds on an open set of sufficiently large
\(\Gamma\), every coefficient in the Laurent expansion in \(1/\Gamma\)
vanishes. Thus all \(M_n=0\), and in particular statement 1 holds.
\(\square\)

\textbf{Important:} This exact theorem is not restricted to the
convergence domain of the Neumann series. The Cayley--Hamilton/adjugate
representation extends it to the entire region in which \(A_\Gamma\) is
invertible.

To claim an exact transfer zero over an entire domain in \(z\), the
equivalent conditions above must hold at every \(z\). Numerical zeros at
finitely many frequencies do not prove an identically vanishing
transfer.

\subsection{5.7 Sector graph and suppression
order}\label{sector-graph-and-suppression-order}

Decompose the coherence space into orthogonal sectors,

\[
\mathcal V=\bigoplus_a\mathcal V_a
\]

and suppose that \(D\) and \(A_{\rm diag}\) preserve each sector, with

\[
A_0=A_{\rm diag}+V
\]

. Regard one action of \(V\) as one edge in the sector graph, and let
\(d\) be the minimum number of \(V\)-insertions required to connect the
source sector to the readout sector. Then

\[
M_n=0\quad(n<d),
\]

\[
\boxed{
K_\Gamma=O(\Gamma^{-1-d})
}
\]

. If several paths undergo phase cancellation, the actual order can be
still higher.

\subsection{5.8 Finite-rate bound}\label{finite-rate-bound}

\[
B_\Gamma=\Gamma D+A_{\rm diag},
\qquad
A_\Gamma=B_\Gamma+V
\]

and assume \(p^\dagger B_\Gamma^{-1}c=0\). The resolvent identity gives

\[
K_\Gamma
=-p^\dagger B_\Gamma^{-1}VA_\Gamma^{-1}c.
\]

\[
r_d=\frac{\|D^{-1}A_{\rm diag}\|}{\Gamma}<1,
\qquad
r=\frac{\|D^{-1}(A_{\rm diag}+V)\|}{\Gamma}<1
\]

then

\[
\boxed{
|K_\Gamma|
\le
\frac{\|p\|\|c\|\|D^{-1}\|^2\|V\|}
{\Gamma^2(1-r_d)(1-r)}
}
\]

.

\subsection{\texorpdfstring{5.9 Precise scope of the condition
\(D>0\)}{5.9 Precise scope of the condition D\textgreater0}}\label{precise-scope-of-the-condition-d0}

The earlier statement that ``the moment theorem requires \(D>0\)'' is
corrected as follows.

\subsubsection{(A) Exact algebraic transfer
zero}\label{a-exact-algebraic-transfer-zero}

Theorem 4 requires only that

\[
\boxed{D\ \text{is invertible}}
\]

. Hermitian positivity is unnecessary.

\subsubsection{(B) Neumann asymptotic
expansion}\label{b-neumann-asymptotic-expansion}

The required conditions are

\[
D\ \text{is invertible},
\qquad
\|D^{-1}A_0\|/\Gamma<1.
\]

. Positivity is unnecessary.

\subsubsection{(C) Physical interpretation as fast damping and uniform
stability}\label{c-physical-interpretation-as-fast-damping-and-uniform-stability}

A clean sufficient condition ensuring that large \(\Gamma\) damps every
mode is

\[
\boxed{
\operatorname{Re}D
=\frac{D+D^\dagger}{2}
\ge d_{\min}I>0
}
\]

. In particular, \(D=D^\dagger>0\) is sufficient. In this case,
\(\Gamma d_{\min}\) provides a coercive damping scale.

\subsubsection{(D) Hermitian metric
overlap}\label{d-hermitian-metric-overlap}

\[
H_A=(A_\Gamma+A_\Gamma^\dagger)/2>0
\]

The bounded geometry proxy \(\eta_H\) defined using

\subsubsection{(E) Singular-damping protected-channel
theorem}\label{e-singular-damping-protected-channel-theorem}

For the present Schur-complement theorem using the orthogonal projector
\(P=\operatorname{Proj}(\ker D)\), assume

\[
\boxed{D=D^\dagger\ge0}
\]

. A non-normal singular \(D\) requires a Riesz projector and a
generalized nullspace, and is outside the scope of the present main
theorem.

\subsubsection{(F) Nonnormal fast
generator}\label{f-nonnormal-fast-generator}

Even when \(D\) is invertible, a non-accretive \(D\) can produce a large
finite-\(\Gamma\) response through pseudospectral amplification.
Algebraic scaling may still be used, but monotonic damping suppression
or a simple visibility bound must not be claimed without checking the
relevant resolvent norm.

\begin{center}\rule{0.5\linewidth}{0.5pt}\end{center}

\section{6. Exact symmetry transfer zero on the full
generator}\label{exact-symmetry-transfer-zero-on-the-full-generator}

\subsection{6.1 Reduced-space projector
theorem}\label{reduced-space-projector-theorem}

Suppose the response space on which \(A(z)\) acts admits orthogonal
projectors

\[
P_\alpha P_\beta=\delta_{\alpha\beta}P_\alpha,
\qquad
\sum_\alpha P_\alpha=I
\]

such that

\[
\boxed{[A(z),P_\alpha]=0\quad(\forall\alpha,z)}
\]

. If the source and readout lie in different sectors,

\[
P_\beta c=c,
\qquad
p^\dagger P_\alpha=p^\dagger,
\qquad
\alpha\neq\beta
\]

then, at every point where \(A(z)\) is invertible,

\[
\boxed{p^\dagger A(z)^{-1}c=0}
\]

.

\textbf{Proof.} If \([A,P_\alpha]=0\), then \([A^{-1},P_\alpha]=0\).
Therefore,

\[
p^\dagger A^{-1}c
=p^\dagger P_\alpha A^{-1}P_\beta c
=p^\dagger A^{-1}P_\alpha P_\beta c=0.
\]

\(\square\)

The earlier formulation in terms of a Hermitian symmetry operator
\(Q=Q^\dagger\) is a special case of this theorem, obtained from its
spectral projectors \(P_\alpha\).

\subsection{6.2 Full-Liouvillian transfer-zero
theorem}\label{full-liouvillian-transfer-zero-theorem}

Let \(\mathcal L_c\) be the zeroth-order Liouvillian including the
control but excluding the probe. Let \(\rho_0\) be the unique steady
state and \(\langle\!\langle\mathbb 1|\) the trace functional, and
define

\[
\Pi_0=|\rho_0\rangle\!\rangle
\langle\!\langle\mathbb 1|,
\qquad
Q_0=I-\Pi_0
\]

.

Suppose the Liouville space admits orthogonal sector projectors
\(\mathbb P_\alpha\) satisfying

\[
[\mathcal L_c,\mathbb P_\alpha]=0,
\qquad
[Q_0,\mathbb P_\alpha]=0
\]

. Define the probe source and detector by

\[
|b_p\rangle\!\rangle
=\mathcal V_p|\rho_0\rangle\!\rangle,
\qquad
\langle\!\langle d_p|
\]

and assume

\[
\mathbb P_\beta|b_p\rangle\!\rangle=|b_p\rangle\!\rangle,
\qquad
\langle\!\langle d_p|\mathbb P_\alpha
=\langle\!\langle d_p|,
\qquad
\alpha\neq\beta
\]

. Then the source-to-detector transfer defined through the group/Drazin
inverse,

\[
\chi_p
=-\langle\!\langle d_p|
\mathcal L_c^\#
|b_p\rangle\!\rangle
\]

vanishes exactly.

\textbf{Proof.} Since \(\mathcal L_c\), \(Q_0\), and
\(\mathbb P_\alpha\) commute, the inverse on \(Q_0\), and hence the
group inverse \(\mathcal L_c^\#\), preserves every sector. The same
orthogonality argument as in the reduced-space theorem makes the
transfer vanish. \(\square\)

\textbf{Interpretive caution.} This theorem directly proves a zero of
the specified source-to-detector transfer. If it is applied to the total
probe susceptibility and yields \(\chi_p=0\), that result may represent
complete transparency and therefore is not, by itself, an EIT no-go. To
obtain an EIT no-go statement, apply the theorem to the targeted sector
correction \(\delta\chi_{\mathbb S}\) defined in Section 10 or to a
cross-sector transfer required to construct it.

\subsection{6.3 Microscopic sufficient condition for Lindblad
symmetry}\label{microscopic-sufficient-condition-for-lindblad-symmetry}

Let \(U\) be a unitary operator on the Hilbert space, and define

\[
\mathscr U(\rho)=U\rho U^\dagger
\]

.

\[
[\mathcal L_c,\mathscr U]=0
\]

The following conditions are sufficient for

\begin{enumerate}
\def\labelenumi{\arabic{enumi}.}
\tightlist
\item
  The full Hamiltonian, including the control, satisfies \[
  U H_cU^\dagger=H_c
  \] .
\item
  The jump operators are closed under unitary mixing: \[
  UL_\mu U^\dagger
  =\sum_\nu V_{\mu\nu}L_\nu,
  \] where \(V\) is unitary on the jump-index space, and the Kossakowski
  matrix/rates are invariant under this transformation.
\item
  If the steady state is unique, then automatically \[
  U\rho_0U^\dagger=\rho_0.
  \] . If the steady state is nonunique, choose a symmetry-invariant
  zeroth-order state.
\end{enumerate}

These conditions include not only Hamiltonian symmetry but also
radiative decay, intersystem crossing, orbital hopping, ground-state
relaxation, dephasing, and correlated jumps.

\subsection{6.4 Full-generator symmetry
audit}\label{full-generator-symmetry-audit}

Before declaring an exact symmetry-protected transfer zero, inspect all
of the following items.

\begin{longtable}[]{@{}
  >{\raggedright\arraybackslash}p{(\columnwidth - 2\tabcolsep) * \real{0.5000}}
  >{\raggedright\arraybackslash}p{(\columnwidth - 2\tabcolsep) * \real{0.5000}}@{}}
\toprule\noalign{}
\begin{minipage}[b]{\linewidth}\raggedright
block
\end{minipage} & \begin{minipage}[b]{\linewidth}\raggedright
required test
\end{minipage} \\
\midrule\noalign{}
\endhead
\bottomrule\noalign{}
\endlastfoot
lower Hamiltonian & Does it preserve the symmetry sector? \\
excited Hamiltonian & Does it preserve the sector after including
spin--orbit, spin--spin, Zeeman, and strain terms? \\
control Hamiltonian & Does it preserve the sector for the selected
polarization and phase? \\
excited-state jumps & Are the jump set and rates symmetry covariant? \\
radiative/ISC loss & Does it avoid mixing the sectors? \\
ground-state dephasing/relaxation & Does it avoid mixing the dark/source
sector? \\
zeroth-order steady state & Is it symmetry invariant? \\
probe source & Does it belong to one well-defined sector? \\
detector & Does it belong to an orthogonal sector different from the
source? \\
\end{longtable}

If any one of these conditions fails, an exact symmetry-protected
transfer zero does not follow. The symmetry-breaking term should then be
treated as \(V\), and the moment hierarchy or perturbative opening
should be evaluated.

\subsection{6.5 Symmetry-protected zero、Krylov zero、accidental
zero}\label{symmetry-protected-zerokrylov-zeroaccidental-zero}

\begin{itemize}
\tightlist
\item
  \textbf{Symmetry-protected transfer zero:} enforced by reducing
  projectors of the full generator and robust against
  symmetry-preserving perturbations.
\item
  \textbf{Krylov exact transfer zero:} every finite transfer path is
  orthogonal to the readout. This can occur even without an explicit
  symmetry.
\item
  \textbf{Accidental parameter zero:} pathway cancellation occurs only
  at a particular detuning or field value; it is not a domain-wide zero.
\end{itemize}

These are called EIT no-go results only when the vanishing transfer is
the targeted sector correction of Section 10 or a cross-channel required
for that correction. They must be reported separately from a zero of the
total susceptibility.

\section{7. Singular Damping and Protected
Channels}\label{singular-damping-and-protected-channels}

The orthogonal protected-subspace theorem in this section assumes

\[
D=D^\dagger\ge0
\]

. If \(D\) is singular, define

\[
P=\operatorname{Proj}(\ker D),
\qquad
Q_D=I-P
\]

. On the fast subspace, define

\[
A_{QQ}
=\Gamma Q_DDQ_D+Q_DA_0Q_D
\]

and introduce the Schur complement for the slow/protected sector,

\[
S_P(\Gamma)
=PA_0P-PA_0Q_DA_{QQ}^{-1}Q_DA_0P
\]

.

If \(PA_0P\) is invertible on the spectral domain of interest,

\[
S_P(\Gamma)=PA_0P+O(\Gamma^{-1})
\]

and

\[
\boxed{
K_\Gamma
=p^\dagger P(PA_0P)^{-1}Pc
+O(\Gamma^{-1})
}
\]

.

Consequently, an \(O(1)\) coherent channel can survive strong damping
only if both the probe and control vectors have projections onto the
genuine kernel of the optical-coherence damping operator.

\begin{center}\rule{0.5\linewidth}{0.5pt}\end{center}

\section{8. Extension to Arbitrary Materials with Multiple Dissipation
Rates}\label{extension-to-arbitrary-materials-with-multiple-dissipation-rates}

\subsection{8.1 Fixed-ratio fast-rate
limit}\label{fixed-ratio-fast-rate-limit}

For a generic material,

\[
A=A_0+\sum_{r=1}^{R}\Gamma_rD_r
\]

. If all fast rates grow along a common scale \(\Gamma\),

\[
\Gamma_r=\Gamma\lambda_r
\]

with fixed \(\lambda_r\), define

\[
D(\boldsymbol\lambda)=\sum_r\lambda_rD_r
\]

. If \(D(\boldsymbol\lambda)\) is invertible on a given ray, the
algebraic moment hierarchy can be applied. If, in addition,

\[
\operatorname{Re}D(\boldsymbol\lambda)>0
\]

, that ray can be interpreted as a uniform fast-damping limit.

If the suppression exponent depends on \(\boldsymbol\lambda\), the
no-go/go classification should be displayed as a phase diagram in
bath-rate space.

\subsection{8.2 Independent-rate
asymptotics}\label{independent-rate-asymptotics}

If the \(\Gamma_r\) vary independently, the leading term can depend on
the path taken toward the asymptotic limit. One should therefore
evaluate the result separately for

\begin{itemize}
\tightlist
\item
  dominant-rate sector
\item
  fixed-ratio rays
\item
  experimentally relevant trajectory \(\Gamma_r(T)\)
\end{itemize}

rather than assigning a single ``universal exponent.''

\subsection{8.3 General finite-dimensional transfer-zero criterion and
singular
points}\label{general-finite-dimensional-transfer-zero-criterion-and-singular-points}

Let \(\theta\) be a parameter vector and consider a finite-dimensional
matrix family \(A(\theta)\). Assume one of the following.

\begin{enumerate}
\def\labelenumi{\arabic{enumi}.}
\tightlist
\item
  \(\theta\) is complex and \(A(\theta)\) is holomorphic on a connected
  complex domain \(\mathcal U\).
\item
  \(\theta\) is real and \(A(\theta)\) is real analytic on a connected
  open real domain \(\mathcal U\).
\end{enumerate}

The identity theorem is not invoked for merely smooth parameter
dependence.

regular set

\[
\mathcal U_{\rm reg}
=\{\theta\in\mathcal U:\det A(\theta)\neq0\}
\]

one has

\[
K(\theta)
=p^\dagger A(\theta)^{-1}c
=
\frac{N(\theta)}{\Delta(\theta)},
\]

\[
N(\theta)=p^\dagger\operatorname{adj}A(\theta)c,
\qquad
\Delta(\theta)=\det A(\theta).
\]

\subsubsection{Theorem 7A --- Regular-domain exact transfer
zero}\label{theorem-7a-regular-domain-exact-transfer-zero}

On a connected analytic domain \(\mathcal U\), if

\[
\boxed{N(\theta)\equiv0}
\]

then, for every \(\theta\in\mathcal U_{\rm reg}\),

\[
\boxed{K(\theta)=0}
\]

. Conversely, if \(K\) vanishes on a nonempty open subset of
\(\mathcal U_{\rm reg}\), then \(N=0\) on that subset, and the
holomorphic or real-analytic identity theorem gives \(N\equiv0\).

\textbf{Caution:} The condition \(\operatorname{adj}A=0\) at a singular
point does not by itself imply that the response is finite or zero. If
the rank deficiency is at least two, the adjugate can vanish
automatically.

\subsubsection{8.3.1 Isolated singularity and Laurent
test}\label{isolated-singularity-and-laurent-test}

Suppose \(\theta_0\) is an isolated singular point and, along a
one-dimensional regularization parameter \(s\),

\[
A(s)^{-1}
=\sum_{j=-\nu}^{\infty}R_j s^j
\]

admits the Laurent expansion above. The transfer is

\[
K(s)=\sum_{j=-\nu}^{\infty}
\left(p^\dagger R_jc\right)s^j.
\]

\begin{itemize}
\tightlist
\item
  The necessary and sufficient condition for pole cancellation is \[
  p^\dagger R_jc=0
  \quad(j=-\nu,\ldots,-1)
  \]
\item
  Laurent finite part： \[
  \boxed{
  \operatorname{FP}_{s=0}K(s)
  =p^\dagger R_0c
  }
  \]
\item
  If every negative-power coefficient vanishes, the scalar-transfer
  singularity is removable, and the analytic continuation is defined by
  \[
  \boxed{
  K_{\rm cont}(0)=p^\dagger R_0c
  }
  \] .
\item
  The continued transfer vanishes at the singular point if all pole
  coefficients vanish and \[
  p^\dagger R_0c=0
  \] .
\end{itemize}

Because the original matrix inverse \(A(0)^{-1}\) does not exist, one
must not write \(K(0)=p^\dagger R_0c\) without explicitly specifying the
regularization.

\subsubsection{8.3.2 Semisimple zero and Drazin/group
inverse}\label{semisimple-zero-and-drazingroup-inverse}

Suppose the zero eigenvalue of \(A\) is semisimple, and let \(P_0\) be
the spectral projector onto its nullspace. Under a shift regularization,

\[
(A+sI)^{-1}
=\frac{P_0}{s}+A^D+O(s),
\]

where \(A^D\) is the Drazin inverse. Hence

\[
K(s)
=\frac{p^\dagger P_0c}{s}
+p^\dagger A^Dc+O(s).
\]

Therefore,

\[
\boxed{p^\dagger P_0c=0}
\]

is the pole-cancellation condition, and

\[
\boxed{
\operatorname{FP}_{s=0}K(s)
=p^\dagger A^Dc
}
\]

. When the pole cancels, this finite part is also the value of the
removable continuation.

For the unique steady-state zero mode of a full Liouvillian, the probe
source is trace zero and therefore normally has no component along the
steady-state projector. The linear response is then defined by the group
inverse on the trace-zero response space.

\subsubsection{8.3.3 Operational rule}\label{operational-rule}

\begin{enumerate}
\def\labelenumi{\arabic{enumi}.}
\tightlist
\item
  At regular points, use the adjugate-numerator identity.
\item
  For a Liouvillian zero mode, use the trace-zero subspace or group
  inverse.
\item
  At an exceptional point or nonsemisimple singularity, inspect every
  Laurent pole coefficient.
\item
  Distinguish the finite part, removable continuation, and ordinary
  inverse.
\item
  Do not treat \(0/0\) as a numerical zero.
\end{enumerate}

\section{9. Weak-probe susceptibility: full Liouvillian derivation and
reduced
kernel}\label{weak-probe-susceptibility-full-liouvillian-derivation-and-reduced-kernel}

\subsection{9.1 Full linear-response
formula}\label{full-linear-response-formula}

Let the probe amplitude be a small complex parameter \(\epsilon\), and
write

\[
\mathcal L(\epsilon)
=\mathcal L_c
+\epsilon\mathcal V_p
+\epsilon^*\widetilde{\mathcal V}_p
+O(|\epsilon|^2)
\]

. Here \(\mathcal L_c\) contains the control, all dissipation, and the
external fields. Let the zeroth-order steady state satisfy

\[
\mathcal L_c\rho^{(0)}=0,
\qquad
\operatorname{Tr}\rho^{(0)}=1
\]

and expand

\[
\rho
=\rho^{(0)}
+\epsilon\rho^{(1)}
+\epsilon^*\widetilde\rho^{(1)}
+O(|\epsilon|^2)
\]

. The coefficient of \(\epsilon\) satisfies

\[
\mathcal L_c\rho^{(1)}
=-\mathcal V_p\rho^{(0)},
\qquad
\operatorname{Tr}\rho^{(1)}=0
\]

.

Assuming a unique stationary state and a finite Liouvillian gap, define
the group inverse \(\mathcal L_c^\#\) on the trace-zero subspace. Then

\[
\boxed{
\rho^{(1)}
=-\mathcal L_c^\#
\mathcal V_p\rho^{(0)}
}
\]

. Apart from the normalization constant \(\chi_0\), the linear
susceptibility of an observable \(O_p\) is

\[
\boxed{
\chi_p
=-\chi_0
\langle\!\langle O_p|
\mathcal L_c^\#
|\mathcal V_p\rho^{(0)}\rangle\!\rangle
}
\]

. This is the reference formula for the full-Liouvillian susceptibility;
the zero eigenvalue of the Liouvillian is not handled with an ordinary
inverse.

\subsection{9.2 Closed two-lower-state first-order
block}\label{closed-two-lower-state-first-order-block}

Among the first-order variables, let \(x\in\mathcal H_e\) be the probe
optical-coherence vector and let the lower-state coherence be

\[
s=\rho_{g_2g_1}^{(1)}
\]

. Assume that the remaining first-order variables either decouple or
have been eliminated exactly by a Schur complement.

Fix the sign convention as

\[
\begin{pmatrix}
A & i\Omega_c d_2/2\\
i\Omega_c^*d_2^\dagger/2 & \gamma_g
\end{pmatrix}
\begin{pmatrix}x\\s\end{pmatrix}
=
\begin{pmatrix}i\Omega_p d_1/2\\0\end{pmatrix}
\]

. Here,

\begin{itemize}
\tightlist
\item
  \(A\): effective optical-coherence generator excluding the control; it
  may include self-energies from other blocks when necessary.
\item
  \(\gamma_g=\gamma_{12}-i\delta_2\): lower-coherence complex
  generator。
\item
  \(d_1,d_2\): probe/control dipole vectors。
\end{itemize}

Assume \(A\) is invertible and set \(G=A^{-1}\). The first row gives

\[
x
=G\frac{i\Omega_p}{2}d_1
-G\frac{i\Omega_c}{2}d_2s.
\]

Substituting this into the second row, without dividing by \(\gamma_g\)
itself, gives

\[
\left(\gamma_g+\beta S_2\right)s
=-\frac{i\Omega_c^*}{2}
 d_2^\dagger G\frac{i\Omega_p}{2}d_1,
\qquad
\beta=\frac{|\Omega_c|^2}{4},
\]

\[
S_2=d_2^\dagger Gd_2
\]

. Therefore, the ideal limit \(\gamma_g=0\) is covered by the same
expression provided that the combined denominator
\(\gamma_g+\beta S_2\neq0\).

\subsection{9.3 Theorem 2A --- Reduced susceptibility as an exact block
Schur
complement}\label{theorem-2a-reduced-susceptibility-as-an-exact-block-schur-complement}

\[
S_i=d_i^\dagger Gd_i,
\qquad
K_{ij}=d_i^\dagger Gd_j
\]

Define the normalized probe response by

\[
\Xi
=\frac{d_1^\dagger x}{i\Omega_p/2}
\]

. Then

\[
\boxed{
\Xi
=S_1-
\frac{\beta K_{12}K_{21}}
{\gamma_g+\beta S_2}
}
\]

.

\textbf{Proof.} The lower-coherence equation above gives

\[
s=
\frac{\Omega_c^*\Omega_p}{4}
\frac{K_{21}}
{\gamma_g+\beta S_2}
\]

. Substituting this into

\[
x
=G\frac{i\Omega_p}{2}d_1
-G\frac{i\Omega_c}{2}d_2s
\]

, acting from the left with \(d_1^\dagger\), and dividing by
\(i\Omega_p/2\) yields the result. This is the Schur complement of the
first-order Liouvillian block and requires no additional adiabatic
approximation. \(\square\)

The control-induced coherent correction is

\[
\boxed{
\Delta\Xi
=\frac{\beta K_{12}K_{21}}
{\gamma_g+\beta S_2}
}
\]

. At a regular point with nonzero control \(\beta\neq0\) and a finite
denominator,

\[
\boxed{K_{12}K_{21}=0}
\]

is the necessary and sufficient zero condition for the complex coherent
correction \(\Delta\Xi\). If the denominator vanishes, use the
Laurent/Drazin analysis of Section 8.3.

\subsection{9.4 Conditions for agreement with the full
Liouvillian}\label{conditions-for-agreement-with-the-full-liouvillian}

The reduced expression above agrees with the full formula under the
following conditions.

\begin{enumerate}
\def\labelenumi{\arabic{enumi}.}
\tightlist
\item
  The response is first order in the probe.
\item
  The zeroth-order state \(\rho^{(0)}\), including the control, has been
  solved correctly.
\item
  The first-order source reduces to the selected probe vector \(d_1\).
\item
  First-order variables other than \((x,s)\) either decouple or are
  absorbed into \(A,\gamma_g,d_i\) through a Schur complement.
\item
  Population saturation and probe-induced redistribution are
  \(O(|\Omega_p|^2)\) and therefore do not enter at first order.
\item
  The detector corresponds to \(d_1^\dagger x\).
\end{enumerate}

If these conditions fail, promote the scalar \(s\) to a ground-coherence
vector and use the full block-matrix Schur complement or
\(\mathcal L_c^\#\) directly.

\subsection{9.5 General matrix ground-coherence
version}\label{general-matrix-ground-coherence-version}

For several ground-state coherences \(s\in\mathbb C^m\), write the
first-order block as

\[
\begin{pmatrix}
A&B\\
C&G_g
\end{pmatrix}
\begin{pmatrix}x\\s\end{pmatrix}
=
\begin{pmatrix}b_p\\0\end{pmatrix}
\]

. Here the source and detector are assumed to belong to the
optical-response sector \(x\).

If \(A\) is invertible, the first row gives

\[
x=A^{-1}b_p-A^{-1}Bs
\]

and, using the ground-coherence Schur complement

\[
\boxed{
S_g=G_g-CA^{-1}B
}
\]

, one obtains

\[
S_gs=-CA^{-1}b_p
\]

. If \(S_g\) is invertible,

\[
\boxed{
s=-S_g^{-1}CA^{-1}b_p
}
\]

and

\[
\boxed{
x=A^{-1}b_p
+A^{-1}BS_g^{-1}CA^{-1}b_p
}
\]

. This form does not require \(G_g\) itself to be invertible and
therefore includes the ideal ground-coherence limit.

Only when \(G_g\) is also invertible may one use the equivalent
alternative Schur complement

\[
\boxed{
x=(A-BG_g^{-1}C)^{-1}b_p
}
\]

.

For an optical detector \(d\),

\[
\chi_{\rm full}
=\chi_0d^\dagger x,
\qquad
\chi_{\rm cut}^{(\mathbb S)}
=\chi_0d^\dagger A^{-1}b_p,
\]

and therefore the response modification mediated by the specified sector
\(\mathbb S\) is

\[
\boxed{
\delta\chi_{\mathbb S}
=\chi_{\rm full}-\chi_{\rm cut}^{(\mathbb S)}
=\chi_0d^\dagger A^{-1}B
S_g^{-1}CA^{-1}b_p
}
\]

. The quantity \(\chi_{\rm cut}^{(\mathbb S)}\) is not the control-off
response; it is an algebraic counterfactual obtained by retaining the
same diagonal blocks, source, detector, and zeroth-order state while
cutting only the targeted \(x\leftrightarrow s\) feedback. The cut
generator need not define an independent CPTP dynamics. Once the sector
\(\mathbb S\) is fixed, we abbreviate it as \(\chi_{\rm cut}\).

If the source or detector also has a direct component in the
ground-coherence sector, define the counterfactual difference using the
full block inverse rather than the expression above. For singular \(A\)
or \(S_g\), use the Laurent/Drazin prescription of Section 8.3.

\subsection{\texorpdfstring{9.6 Recovery of the textbook three-level
\(\Lambda\)
model}{9.6 Recovery of the textbook three-level \textbackslash Lambda model}}\label{recovery-of-the-textbook-three-level-lambda-model}

Consider the standard \(\Lambda\) system with one excited state
\(|e\rangle\) and two lower states \(|g_1\rangle,|g_2\rangle\). Choose
the dipole normalization

\[
d_1=d_2=1
\]

and set

\[
A=\gamma_{31}-i\Delta_p,
\qquad
\gamma_g=\gamma_{21}-i\delta_2,
\qquad
\beta=\frac{|\Omega_c|^2}{4}
\]

. Then

\[
S_1=S_2=K_{12}=K_{21}=A^{-1}.
\]

Substitution into Theorem 2A gives

\[
\begin{aligned}
\Xi
&=\frac1A-
\frac{\beta/A^2}{\gamma_g+\beta/A}\\
&=\boxed{
\frac{\gamma_g}{A\gamma_g+\beta}
}.
\end{aligned}
\]

Therefore, the probe optical coherence is

\[
\boxed{
\rho_{eg_1}^{(1)}
=\frac{i\Omega_p}{2}
\frac{\gamma_{21}-i\delta_2}
{(\gamma_{31}-i\Delta_p)(\gamma_{21}-i\delta_2)+|\Omega_c|^2/4}
}
\]

. This is the standard three-level EIT susceptibility in the adopted
rotating-frame sign convention.

\begin{itemize}
\tightlist
\item
  control off \((\beta=0)\): \(\Xi=A^{-1}\)
\item
  exact two-photon resonance and ideal ground coherence
  \((\delta_2=0,\gamma_{21}\to0)\): \(\Xi\to0\)
\item
  in the same limit,
  \(\delta\Xi_{\mathbb S}=\Xi_{\rm full}-\Xi_{\rm cut}=-A^{-1}\neq0\)
\end{itemize}

Thus, a zero of the total susceptibility represents perfect EIT and
differs from a zero of the sector-mediated correction. This recovery
shows that the framework contains the textbook \(\Lambda\) response as a
special case.

\subsection{9.7 Recovery of a Mishina-type multi-excited-state
model}\label{recovery-of-a-mishina-type-multi-excited-state-model}

Consider an operative excited manifold consisting of \(N_e\) resolved
branches, with an optical-coherence generator diagonal in the selected
basis. Let \(\lambda\) label an inhomogeneous class, velocity class, or
static detuning class, and write

\[
A_\lambda
=\operatorname{diag}
\left(a_1(\lambda),\ldots,a_{N_e}(\lambda)\right),
\qquad
a_j(\lambda)=\gamma_j-i\Delta_j(\lambda)
\]

. For the probe and control dipole vectors

\[
d_1=(d_{11},\ldots,d_{1N_e})^{\mathsf T},
\qquad
d_2=(d_{21},\ldots,d_{2N_e})^{\mathsf T}
\]

, one obtains

\[
\boxed{
S_{1,\lambda}
=\sum_{j=1}^{N_e}
\frac{|d_{1j}|^2}{a_j(\lambda)}
},
\qquad
\boxed{
S_{2,\lambda}
=\sum_{j=1}^{N_e}
\frac{|d_{2j}|^2}{a_j(\lambda)}
}
\]

and

\[
\boxed{
K_{12,\lambda}
=\sum_{j=1}^{N_e}
\frac{d_{1j}^*d_{2j}}{a_j(\lambda)}
},
\qquad
\boxed{
K_{21,\lambda}
=\sum_{j=1}^{N_e}
\frac{d_{2j}^*d_{1j}}{a_j(\lambda)}
}
\]

. Hence the response of each class is

\[
\boxed{
\Xi_\lambda
=S_{1,\lambda}
-
\frac{\beta K_{12,\lambda}K_{21,\lambda}}
{\gamma_g+\beta S_{2,\lambda}}
}.
\]

This is the diagonal-resolvent structure of a Mishina-type multilevel
susceptibility, containing the coherent sum of Raman amplitudes through
multiple excited levels and the self-energy of each optical leg. Static
Doppler or inhomogeneous broadening is incorporated only after the local
response has been calculated, by averaging

\[
\boxed{
\overline\chi
=\int d\lambda\,P(\lambda)\chi_\lambda
}
\]

. When optical pumping is included, the zeroth-order populations enter
\(\rho^{(0)}\) and the source \(b_p\), so they must first be solved with
the full Liouvillian. Thus, the present theory contains the
\textbf{multi-excited-state Green-function interference structure} of
Mishina-type theories; it does not unconditionally identify its
approximations with the propagation, velocity-changing collisions, or
pumping approximations of any particular experiment.

\subsection{9.8 Observable dictionary}\label{observable-dictionary}

\begin{longtable}[]{@{}
  >{\raggedright\arraybackslash}p{(\columnwidth - 4\tabcolsep) * \real{0.3333}}
  >{\raggedright\arraybackslash}p{(\columnwidth - 4\tabcolsep) * \real{0.3333}}
  >{\raggedright\arraybackslash}p{(\columnwidth - 4\tabcolsep) * \real{0.3333}}@{}}
\toprule\noalign{}
\begin{minipage}[b]{\linewidth}\raggedright
object
\end{minipage} & \begin{minipage}[b]{\linewidth}\raggedright
definition
\end{minipage} & \begin{minipage}[b]{\linewidth}\raggedright
role
\end{minipage} \\
\midrule\noalign{}
\endhead
\bottomrule\noalign{}
\endlastfoot
\(K_{12},K_{21}\) & off-diagonal resolvent & coherent Raman/interference
vertex \\
\(S_1,S_2\) & diagonal resolvent & optical response/self-energy \\
\(\chi\) & \(\chi_0\Xi\) & local linear susceptibility \\
\(C_{\rm abs}\) & normalized absorption contrast & specified
observable \\
\(T\) & \(\exp[-L\alpha]\) & propagated transmission \\
\end{longtable}

For a non-normal generator, in general,

\[
K_{21}\neq K_{12}^*.
\]

so the two quantities must be computed independently.

Hermitian part

\[
H_A=(A+A^\dagger)/2>0
\]

Only when

\[
\eta_H=
\frac{|d_1^\dagger H_A^{-1}d_2|^2}
{(d_1^\dagger H_A^{-1}d_1)
(d_2^\dagger H_A^{-1}d_2)}
\]

does the bounded geometry proxy satisfy \(0\le\eta_H\le1\).

\subsection{9.9 Relative absorption contrast and visibility
branches}\label{relative-absorption-contrast-and-visibility-branches}

As one example, define

\[
C_{\rm abs}
=\frac{\operatorname{Im}S_1-
\operatorname{Im}\Xi}
{\operatorname{Im}S_1}
=\frac{\operatorname{Im}\Delta\Xi}
{\operatorname{Im}S_1}
\]

. The observable and sign convention must be fixed consistently in
material-to-material comparisons.

For a one-insertion channel, suppose

\[
K_{12},K_{21}=O(\Gamma^{-2}),
\qquad
S_1,S_2=O(\Gamma^{-1})
\]

and define

\[
y=\left|\frac{\beta S_2}{\gamma_g}\right|
\]

.

\begin{itemize}
\tightlist
\item
  \(y\ll1\): \[
  \Delta\Xi=O(\Gamma^{-4}),
  \qquad
  C_{\rm abs}=O(\Gamma^{-3}).
  \]
\item
  \(y\gg1\): \[
  \Delta\Xi=O(\Gamma^{-3}),
  \qquad
  C_{\rm abs}=O(\Gamma^{-2}).
  \]
\end{itemize}

For fixed \(\gamma_g\neq0\) and fixed \(\Omega_c\), one ultimately has
\(y\to0\), so the strict \(\Gamma\to\infty\) limit is the former branch.

\begin{center}\rule{0.5\linewidth}{0.5pt}\end{center}

\section{10. Full-Liouvillian EIT no-go and observable
criterion}\label{full-liouvillian-eit-no-go-and-observable-criterion}

\subsection{10.1 Pathway-resolved response
operator}\label{pathway-resolved-response-operator}

For a weak-probe problem with a unique steady state, define on the
trace-zero response space

\[
\mathbb M=-Q_0\mathcal L_cQ_0
\]

as the response generator, so that

\[
\mathbb M|u^{(1)}\rangle\!\rangle
=|b_p\rangle\!\rangle
\]

. Let \(\mathbb S\) be the long-lived coherence sector responsible for
the target EIT mechanism, and let \(\mathbb X\) be the complementary
direct optical-response sector. Partition the response problem as

\[
\mathbb M=
\begin{pmatrix}
A&B\\
C&G_g
\end{pmatrix},
\qquad
|b_p\rangle\!\rangle=
\begin{pmatrix}b_p\\0\end{pmatrix},
\qquad
\langle\!\langle d_p|=
\begin{pmatrix}d^\dagger&0\end{pmatrix}
\]

.

Define the full response and the counterfactual response obtained by
cutting only the targeted cross-feedback as

\[
\chi_{\rm full}
=\chi_0d^\dagger x,
\qquad
\chi_{\rm cut}^{(\mathbb S)}
=\chi_0d^\dagger A^{-1}b_p
\]

. The cut operation sets \(B=C=0\) while holding \(A,G_g,b_p,d\) and the
zeroth-order state fixed. It is therefore generally different from
switching the control field off.

\subsection{10.2 Theorem 2B --- Targeted full-Liouvillian
sector-resolved EIT
no-go}\label{theorem-2b-targeted-full-liouvillian-sector-resolved-eit-no-go}

If \(A\) and

\[
S_g=G_g-CA^{-1}B
\]

are invertible at a regular point,

\[
\boxed{
\delta\chi_{\mathbb S}
\equiv
\chi_{\rm full}-\chi_{\rm cut}^{(\mathbb S)}
=
\chi_0d^\dagger A^{-1}B
S_g^{-1}CA^{-1}b_p
}
\]

.

On a specified connected analytic parameter domain \(\mathcal U\), if

\[
\boxed{
\delta\chi_{\mathbb S}(\theta)\equiv0
\quad(\theta\in\mathcal U_{\rm reg})
}
\]

then the net \(\mathbb S\)-mediated contribution to the specified
source-to-detector response vanishes, and the configuration is an
\textbf{exact \(\mathbb S\)-sector EIT no-go}. This does not require the
internal block couplings \(B\) and \(C\) to vanish individually;
complete cancellation at the detector is also included.

Conversely, \(\delta\chi_{\mathbb S}\neq0\) establishes permission for a
net pathway through the specified sector, but it is not sufficient for a
narrow transparency feature, positive contrast, or experimental
visibility. Only when \(\mathbb S\) is a long-lived coherence sector and
\(B,C\) originate from the control Hamiltonian is this called coherent
EIT-pathway permission.

For a closed scalar block with two lower states,

\[
\delta\Xi_{\mathbb S}
=\Xi_{\rm full}-\Xi_{\rm cut}
=-\frac{\beta K_{12}K_{21}}
{\gamma_g+\beta S_2}.
\]

At a regular point with nonzero control and a finite denominator,

\[
\boxed{
\delta\Xi_{\mathbb S}=0
\Longleftrightarrow
K_{12}K_{21}=0
}
\]

.

\textbf{Structural zeros versus pointwise zeros.} A value
\(\delta\chi_{\mathbb S}=0\) at one detuning or field value is not a
domain-wide no-go. An exact no-go must be established as an identity
using a reducing symmetry, finite Krylov orthogonality, an
adjugate-numerator identity, or an analytic identity theorem.

The quantity \(\delta\chi_{\mathbb S}\) is the difference between the
full response and the \(\mathbb S\)-cut response and is therefore a
finite-dimensional rational transfer. When using a Krylov/moment
theorem, apply it either to required cross-transfers such as
\(d^\dagger A^{-1}B\) and \(CA^{-1}b_p\), or construct an augmented
realization that represents the full/cut difference as a single
transfer.

\subsection{10.3 Sector dependence, invariance, and
interpretation}\label{sector-dependence-invariance-and-interpretation}

\subsubsection{\texorpdfstring{10.3.1 Physical dependence on
\(\mathbb S\)}{10.3.1 Physical dependence on \textbackslash mathbb S}}\label{physical-dependence-on-mathbb-s}

The quantity \(\delta\chi_{\mathbb S}\) depends on the specified
projector pair

\[
P_X+P_S=I,
\qquad
P_XP_S=0,
\qquad
\operatorname{Ran}P_S=\mathbb S
\]

. In projector form,

\[
\mathbb M_{\rm cut}^{(\mathbb S)}
=P_X\mathbb MP_X+P_S\mathbb MP_S
\]

and

\[
\boxed{
\delta\chi_{\mathbb S}
=\chi_0\langle\!\langle d_p|
\left(\mathbb M^{-1}-(\mathbb M_{\rm cut}^{(\mathbb S)})^{-1}\right)
|b_p\rangle\!\rangle
}
\]

. Therefore, every no-go statement is a statement about

\[
(\mathfrak C,\mathbb S,\mathbb X,b_p,d_p)
\]

and is not a system-wide scalar label independent of the sector choice.

\subsubsection{10.3.2 Block-preserving basis
invariance}\label{block-preserving-basis-invariance}

Let \(U=U_X\oplus U_S\) be an invertible basis transformation that does
not mix \(\mathbb X\) and \(\mathbb S\). If the generator, source, and
detector are transformed consistently,

\[
\boxed{
\delta\chi_{\mathbb S}'=\delta\chi_{\mathbb S}
}
\]

. Hence the result is independent of the basis chosen within either
sector.

By contrast, if one performs a transformation that mixes \(\mathbb X\)
and \(\mathbb S\) and then redefines the same coordinate cut, one has
generally selected a different projector, and \(\delta\chi\) changes.
This is not basis dependence; it is a \textbf{change in the physical
pathway definition}.

\subsubsection{10.3.3 Algebraic
counterfactual}\label{algebraic-counterfactual}

The operator \(\mathbb M_{\rm cut}^{(\mathbb S)}\) is an algebraic
counterfactual used for pathway attribution and need not itself be an
independent trace-preserving completely positive generator. Accordingly,

\begin{itemize}
\tightlist
\item
  do not identify \(\chi_{\rm cut}^{(\mathbb S)}\) with the control-off
  spectrum;
\item
  do not identify \(\delta\chi_{\mathbb S}\) unconditionally with a
  directly measurable observable;
\item
  for experimental comparison, use \(\chi_{\rm full}\), control scans,
  two-photon dependence, ground-coherence dependence, and pole/residue
  analysis.
\end{itemize}

\subsubsection{10.3.4 Additional conditions for a coherent
interpretation}\label{additional-conditions-for-a-coherent-interpretation}

Generic blocks \(B,C\) can contain dissipative transfer as well as
Hamiltonian coupling. The quantity \(\delta\chi_{\mathbb S}\) is called
a coherent EIT correction only if the following conditions hold.

\begin{enumerate}
\def\labelenumi{\arabic{enumi}.}
\tightlist
\item
  \(\mathbb S\) is a long-lived off-diagonal density-matrix sector
  between the selected lower states.
\item
  The EIT-relevant leading terms in \(B,C\) originate from the control
  Hamiltonian.
\item
  Population pumping and dissipative cross-transfer have been separated
  and audited in the full Liouvillian.
\end{enumerate}

\subsection{10.4 Total susceptibility zero is not an EIT
no-go}\label{total-susceptibility-zero-is-not-an-eit-no-go}

For an ideal scalar \(\Lambda\) block, take

\[
A=1,
\qquad
d_1=d_2=1,
\qquad
\beta=1,
\qquad
\gamma_g=0
\]

. Then

\[
S_1=S_2=K_{12}=K_{21}=1,
\]

and hence

\[
\Xi_{\rm cut}=1,
\qquad
\Xi_{\rm full}
=1-\frac{1}{1}=0.
\]

On the other hand, define the transparency correction by

\[
\Delta\Xi_{\rm tr}
\equiv
\Xi_{\rm cut}-\Xi_{\rm full}
\]

. Then

\[
\Delta\Xi_{\rm tr}=1\neq0.
\]

Therefore,

\[
\boxed{
\Xi_{\rm full}=0
}
\]

can represent perfect transparency. It must not be defined as a no-go
condition.

\subsection{10.5 General projector and singular
formulation}\label{general-projector-and-singular-formulation}

If explicit \((x,s)\) coordinates cannot be chosen, use response-space
projectors \(P_X,P_S\) and define

\[
\mathbb M_{\rm cut}^{(\mathbb S)}
=P_X\mathbb MP_X+P_S\mathbb MP_S
\]

. When the source and detector belong to the \(X\) sector,

\[
\chi_{\rm full}
=\chi_0\langle\!\langle d_p|
\mathbb M^{-1}
|b_p\rangle\!\rangle,
\]

\[
\chi_{\rm cut}^{(\mathbb S)}
=\chi_0\langle\!\langle d_p|
(\mathbb M_{\rm cut}^{(\mathbb S)})^{-1}
|b_p\rangle\!\rangle.
\]

If zero modes or isolated singularities occur, do not use an ordinary
inverse. Use a group/Drazin inverse on the compatible trace-zero
subspace, or a Laurent finite part/removable continuation.

The partition \(X\oplus S\) must be physically specified. Choosing a
different coherence sector means classifying a different targeted EIT
channel.

\subsection{10.6 Observable EIT
criterion}\label{observable-eit-criterion}

Even when the configuration is not an exact \(\mathbb S\)-sector no-go,
at least the following must be verified before claiming observable EIT.

\begin{enumerate}
\def\labelenumi{\arabic{enumi}.}
\tightlist
\item
  \(\delta\chi_{\mathbb S}\neq0\) for the physically specified sector
  \(\mathbb S\).
\item
  There is an absorption reduction, or a corresponding transmission
  increase, near two-photon resonance.
\item
  The feature depends on the ground-coherence decay rate.
\item
  Control-power scans, pole/residue analysis, linewidth scaling, or
  equivalent tests show that the feature cannot be explained solely by
  Autler--Townes splitting.
\item
  The signal exceeds the threshold after including ensemble averaging,
  propagation, background, and signal-to-noise ratio.
\end{enumerate}

A control-on/off difference alone also contains population pumping,
power broadening, and saturation. Use the pathway cut together with
\(\gamma_g\) dependence, two-photon dependence, and pole tracking to
isolate the coherent EIT component.

\subsection{10.7 Validation tests}\label{validation-tests}

At minimum, compare the following on the same parameter grid.

\begin{enumerate}
\def\labelenumi{\arabic{enumi}.}
\tightlist
\item
  A trace-zero constrained solve of the full GKSL Liouvillian and a
  group/Drazin solve
\item
  The full first-order Liouvillian response and the exact Schur-reduced
  formula
\item
  \(\chi_{\rm full}\)、\(\chi_{\rm cut}^{(\mathbb S)}\)、\(\delta\chi_{\mathbb S}\)
\item
  In the textbook \(\Lambda\) limit, \[
  \chi_{\rm full}=0,
  \qquad
  \delta\chi_{\mathbb S}\neq0
  \]
\item
  Agreement between the Mishina-type sum formula and direct block
  inversion for a diagonal multi-excited-state block
\item
  Validity of the \(S_g\) formulation even when \(G_g\) is singular
\item
  Invariance of \(\delta\chi_{\mathbb S}\) under block-preserving basis
  transformations
\item
  The generic change in \(\delta\chi\) when the projector is redefined
  after an \(X\)-\(S\) mixing transformation
\end{enumerate}

The numerical tests in Section 21.3 are fixed as core solver regression
tests.

\section{11. Symmetry and group-theoretical
criterion}\label{symmetry-and-group-theoretical-criterion}

\subsection{11.1 Leading moment as projected dipole
overlap}\label{leading-moment-as-projected-dipole-overlap}

If \(D\propto I\), or is scalar within the sector,

\[
M_0\propto p^\dagger c
\]

. For a complete operative manifold \(E\),

\[
p^\dagger c
=\langle g_1|
\hat d_p^\dagger\Pi_E\hat d_c
|g_2\rangle.
\]

Thus, the leading-order go/no-go result is determined by the
lower-manifold matrix element of the effective tensor operator

\[
\hat O_{pc}=\hat d_p^\dagger\Pi_E\hat d_c
\]

.

\subsection{11.2 Point-group selection
rule}\label{point-group-selection-rule}

A nonzero matrix element requires

\[
\Gamma(g_1)^*\otimes
\Gamma(\hat O_{pc})\otimes
\Gamma(g_2)
\]

to contain the totally symmetric representation. The representation
\(\Gamma(\hat O_{pc})\) is selected by the polarization tensor
\(\epsilon_p^*\otimes\epsilon_c\), the operative manifold, and any
cavity projection.

\subsection{11.3 Spin-scalar dipole}\label{spin-scalar-dipole}

If the electric dipole operator is spin scalar to leading order,

\[
\hat{\mathbf d}=\hat{\mathbf d}_{\rm orb}\otimes I_{\rm spin}
\]

. For a lower-state pair that shares the same orbital state and differs
only by orthogonal spin states, the scalar leading moment tends to
vanish.

However, a finite higher-order moment can open in the presence of

\begin{itemize}
\tightlist
\item
  transverse spin-orbit coupling
\item
  nonsecular spin-spin coupling
\item
  hyperfine mixing
\item
  transverse magnetic field
\item
  spin-dependent electric-dipole correction
\item
  incomplete manifold selection
\item
  branch-dependent damping
\end{itemize}

Therefore, an exact no-go and asymptotic suppression must be
distinguished in a physical system.

\subsection{11.4 Caution concerning orbital
degeneracy}\label{caution-concerning-orbital-degeneracy}

An orbital-degenerate lower manifold is not sufficient for a go result.
The following are required.

\begin{enumerate}
\def\labelenumi{\arabic{enumi}.}
\tightlist
\item
  The lower-state symmetry permits the transition.
\item
  The polarization tensor contains the required irreducible
  representation.
\item
  The reduced dipole matrix element is nonzero.
\item
  The relevant spin overlap is nonzero.
\item
  The operative excited manifold has been selected correctly.
\item
  Damping and ground-state coherence permit practical visibility.
\end{enumerate}

.

\begin{center}\rule{0.5\linewidth}{0.5pt}\end{center}

\section{12. Classification taxonomy}\label{classification-taxonomy}

Classify an arbitrary configuration \(\mathfrak C\) as follows.

\subsection{Class 0 --- Unspecified /
underdetermined}\label{class-0-unspecified-underdetermined}

The required Hamiltonian, dipoles, jump operators, rates, or observable
are insufficient, so no classification can be made.

\subsection{Class I --- Optical dark-subspace
capable}\label{class-i-optical-dark-subspace-capable}

\[
\dim\ker\Omega>0.
\]

This establishes only the existence of an instantaneous optical dark
superposition. A stationary dark state additionally requires the full
Lindblad conditions of Section 3.4, and observable EIT is not
guaranteed.

\subsection{Class II --- Exact sector-resolved EIT
no-go}\label{class-ii-exact-sector-resolved-eit-no-go}

On the target domain,

\[
\boxed{
\delta\chi_{\mathbb S}(\theta)\equiv0
}
\]

is proved by the finite Krylov criterion, an adjugate identity, or a
reducing symmetry of the full generator that separates the targeted
cross-transfer.

For a regular closed block with two lower states, this is equivalent to

\[
K_{12}K_{21}\equiv0
\]

. A zero of the total susceptibility, \(\chi_{\rm full}=0\), is not the
definition of this class.

\subsection{Class III --- Asymptotic sector-mediated
suppression}\label{class-iii-asymptotic-sector-mediated-suppression}

Let the first nonzero moments of the relevant cross-transfers be
\(m_{12}\) and \(m_{21}\), respectively, so that

\[
K_{12}=O(\Gamma^{-m_{12}-1}),
\qquad
K_{21}=O(\Gamma^{-m_{21}-1}).
\]

At finite \(\Gamma\) the response is not exactly zero, but fast damping
suppresses the targeted \(\mathbb S\)-sector correction systematically.
For example, in a scalar block with \(S_2=O(\Gamma^{-1})\),

\[
\delta\Xi_{\mathbb S}=
\begin{cases}
O(\Gamma^{-m_{12}-m_{21}-2}), & |\beta S_2/\gamma_g|\ll1,\\
O(\Gamma^{-m_{12}-m_{21}-1}), & |\beta S_2/\gamma_g|\gg1.
\end{cases}
\]

When classifying only one transfer, record \(K=O(\Gamma^{-m-1})\), where
\(m\) is its first nonzero moment.

\subsection{Class IV --- Leading-order sector-kernel
go}\label{class-iv-leading-order-sector-kernel-go}

In the regular full-rank damping limit,

\[
K_{12}=O(\Gamma^{-1}),
\qquad
K_{21}=O(\Gamma^{-1})
\]

and the product of the leading coefficients is nonzero, the targeted
\(\mathbb S\)-sector pathway is permitted at leading order. It is
interpreted as a coherent pathway when \(\mathbb S\) is a long-lived
coherence and the cross-coupling is Hamiltonian in origin. This is still
not sufficient for observable EIT.

\subsection{Class V --- Protected go}\label{class-v-protected-go}

Under singular damping, if the leading protected-subspace transfer
satisfies

\[
\boxed{
p^\dagger P(PA_0P)^{-1}Pc\neq0
}
\]

then an \(O(1)\) channel remains. Nonzero source and detector
projections alone do not exclude phase cancellation.

\subsection{Class VI --- Practical
no-go}\label{class-vi-practical-no-go}

\[
K_{12}K_{21}\neq0
\]

Even if the kernel is nonzero, the full observable response may satisfy

\[
|C_{\rm EIT}|<C_{\rm min}
\]

or fail the constraints imposed by signal-to-noise ratio, laser power,
ionization, heating, or optical depth.

\subsection{Class VII --- Observable go}\label{class-vii-observable-go}

A narrow coherent transparency feature identifiable as EIT is obtained
for the specified observable, normalization, and detection threshold.

\subsection{Class VIII --- Tunable
transition}\label{class-viii-tunable-transition}

Fields, strain, polarization, cavities, or temperature move the
configuration among Classes II--VII.

\begin{center}\rule{0.5\linewidth}{0.5pt}\end{center}

\section{13. Standard Procedure for Applying the Theory to an Arbitrary
Material}\label{standard-procedure-for-applying-the-theory-to-an-arbitrary-material}

\subsection{Step 1 --- Fix the
configuration}\label{step-1-fix-the-configuration}

Specify the following uniquely.

\begin{itemize}
\tightlist
\item
  lower-state pair or lower manifold
\item
  operative excited manifold
\item
  probe/control frequencies and polarizations
\item
  field、strain、cavity orientation
\item
  temperature and bath condition
\item
  measured observable
\item
  targeted long-lived coherence sector \(\mathbb S\)
\end{itemize}

\subsection{Step 2 --- Construct the microscopic
input}\label{step-2-construct-the-microscopic-input}

The required inputs are listed below.

\begin{longtable}[]{@{}
  >{\raggedright\arraybackslash}p{(\columnwidth - 2\tabcolsep) * \real{0.5000}}
  >{\raggedright\arraybackslash}p{(\columnwidth - 2\tabcolsep) * \real{0.5000}}@{}}
\toprule\noalign{}
\begin{minipage}[b]{\linewidth}\raggedright
category
\end{minipage} & \begin{minipage}[b]{\linewidth}\raggedright
required data
\end{minipage} \\
\midrule\noalign{}
\endhead
\bottomrule\noalign{}
\endlastfoot
coherent structure & \(H_g,H_e\)、fine
structure、Zeeman、Stark、strain、hyperfine \\
optical coupling & complex dipole matrix elements、selection
rules、polarization geometry \\
excited dissipation & radiative decay、ISC、pure dephasing、orbital/spin
relaxation \\
ground dissipation & \(T_1,T_2,T_2^*\)、population exchange \\
disorder & strain distribution、spectral diffusion、orientation
distribution \\
experimental & \(\Omega_c\)、optical depth、sample length、detection
noise \\
\end{longtable}

Retain the complex phases, not only the absolute oscillator strengths.

\subsection{Step 3 --- Define the full master
equation}\label{step-3-define-the-full-master-equation}

\[
\dot\rho=\mathcal L(\rho)
\]

Construct the equation and distinguish, for every jump operator,

\begin{itemize}
\tightlist
\item
  jumps within a manifold
\item
  loss out of the manifold
\item
  correlated jump
\item
  quasi-static disorder
\end{itemize}

.

\subsection{Step 4 --- Determine whether a reduced kernel is
valid}\label{step-4-determine-whether-a-reduced-kernel-is-valid}

If the fixed-ground optical-coherence block is invariant, use
\(A_j^{-1}\). Otherwise, use the full Liouville-space resolvent.

\subsection{Step 5 --- Compute the dark-state
rank}\label{step-5-compute-the-dark-state-rank}

\[
\dim\mathcal D_{\rm dark}
=N_g-\operatorname{rank}\Omega.
\]

\subsection{Step 6 --- Compute the matrix
kernel}\label{step-6-compute-the-matrix-kernel}

\[
\mathbf K(z,T,\mathbf F)
=C^\dagger A^{-1}C
\]

or its Liouville-space counterpart.

At minimum, save

\[
S_1,S_2,K_{12},K_{21}
\]

on the same grid.

\subsection{Step 7 --- Test the targeted sector
zero}\label{step-7-test-the-targeted-sector-zero}

Define the target sector \(\mathbb S\), and construct
\(\chi_{\rm cut}^{(\mathbb S)}\) and \(\delta\chi_{\mathbb S}\). Test an
exact no-go in the following order.

\begin{enumerate}
\def\labelenumi{\arabic{enumi}.}
\tightlist
\item
  A reducing symmetry that separates the targeted cross-transfer
\item
  finite Krylov moments
\item
  adjugate numerator identity
\item
  symbolic group-theory selection rule
\item
  A numerical zero scan only as supporting evidence
\end{enumerate}

Do not use a zero of the total susceptibility as the no-go criterion.

\subsection{Step 8 --- Determine the strong-damping
order}\label{step-8-determine-the-strong-damping-order}

\[
M_0,M_1,\ldots
\]

Compute the moments and identify the first nonzero one. For multiple
rates, use the experimentally relevant trajectory \(\Gamma_r(T)\).

\subsection{Step 9 --- Compute the full response and the sector-resolved
correction}\label{step-9-compute-the-full-response-and-the-sector-resolved-correction}

At minimum, compute

\[
\chi_{\rm full},
\quad
\chi_{\rm cut}^{(\mathbb S)},
\quad
\delta\chi_{\mathbb S},
\quad
C_{\rm abs},
\quad
T
\]

and convert them to the experimental observable, such as fluorescence,
when necessary.

\subsection{Step 10 --- Distinguish EIT from
ATS}\label{step-10-distinguish-eit-from-ats}

Do not identify EIT from a transparency dip alone. At minimum, inspect

\begin{itemize}
\tightlist
\item
  ground-coherence dependence
\item
  control-power dependence
\item
  pole/eigenmode structure
\item
  two-photon linewidth
\item
  the difference from the pathway-cut response
\item
  the difference from the control-off response
\end{itemize}

. Autler--Townes splitting generated by a strong control must be
reported separately from coherent dark-state EIT.

\subsection{Step 11 --- Perform ensemble averaging
correctly}\label{step-11-perform-ensemble-averaging-correctly}

For a quasi-static parameter \(\xi\),

\[
\bar\chi(\omega)
=\int d\xi\,P(\xi)\chi(\omega;\xi)
\]

and average the observable response. A simple substitution of an average
Hamiltonian or average linewidth can be invalid.

\subsection{Step 12 --- Report the classification and
uncertainty}\label{step-12-report-the-classification-and-uncertainty}

At minimum, output

\[
\{\text{class},m,K_{12},K_{21},\delta\chi_{\mathbb S},C_{\rm EIT},\text{threshold margin},\text{uncertainty}\}
\]

.

\begin{center}\rule{0.5\linewidth}{0.5pt}\end{center}

\section{14. Separation of the Markovian Main Scope, Static Broadening,
and Out-of-Scope
Extensions}\label{separation-of-the-markovian-main-scope-static-broadening-and-out-of-scope-extensions}

\subsection{14.1 Fast Markovian dynamics}\label{fast-markovian-dynamics}

When the system has the finite-dimensional time-independent GKSL
generator and finite spectral gap assumed in the main theory, use
resolvent, moment, and Schur-complement analyses.

\subsection{14.2 Quasi-static disorder}\label{quasi-static-disorder}

Slow strain fluctuations, orientation distributions, and static spectral
diffusion are not ergodic fast damping. Compute the response for each
configuration and then perform the ensemble average.

Even for the same total linewidth,

\begin{itemize}
\tightlist
\item
  fast homogeneous exchange
\item
  static inhomogeneous broadening
\end{itemize}

can exhibit different EIT survival.

\subsection{14.3 Frequency-dependent self-energy --- outside the main
theorem}\label{frequency-dependent-self-energy-outside-the-main-theorem}

A genuinely non-Markovian bath lies outside the main theorem of this
version. For formal comparison, write

\[
A(z)=z-H_e-\Sigma(z)
\]

and retain the transfer-kernel definition

\[
K(z)=p^\dagger A(z)^{-1}c
\]

.

However, to use the finite-moment theorem, one needs, for example, on
the target band,

\[
\Sigma(z)=\Gamma D+\Sigma_0(z)
\]

and

\[
\operatorname{Re}D\ge d_{\min}I>0,
\qquad
\|D^{-1}\Sigma_0(z)\|/\Gamma<1
\]

or another uniformly accretive factorization. In a Hermitian model,
\(D>0\) is a special case.

Do not apply the finite-dimensional Cayley--Hamilton criterion without
modification to an arbitrary non-Markovian bath, branch cut, or
continuum.

\subsection{14.4 Continuum excited manifolds --- controlled
approximations
only}\label{continuum-excited-manifolds-controlled-approximations-only}

For a continuum or a very large vibronic manifold, use

\begin{itemize}
\tightlist
\item
  controlled truncation
\item
  tensor-network/open-system solver
\item
  spectral-density integral equation
\item
  operator-valued transfer function
\end{itemize}

. Apply the finite-dimensional classification only as a controlled
approximation after demonstrating convergence with respect to
truncation.

\begin{center}\rule{0.5\linewidth}{0.5pt}\end{center}

\section{15. Lifting a No-Go Condition by External
Control}\label{lifting-a-no-go-condition-by-external-control}

\subsection{15.1 Generic perturbation}\label{generic-perturbation}

Add a sector-changing perturbation

\[
V_{\rm mix}=\lambda V_1+O(\lambda^2)
\]

to a system with an exact or leading-order no-go condition.

If one insertion opens a path between the source and readout sectors and
the reciprocal kernels \(K_{12}\) and \(K_{21}\) open at the same order,

\[
K_{12},K_{21}=O(\lambda\Gamma^{-2}),
\qquad
K_{12}K_{21}=O(\lambda^2\Gamma^{-4}).
\]

For a non-normal or nonreciprocal model, determine the orders of the two
kernels separately.

The canonical relative absorption contrast therefore scales as

\[
C_{\rm abs}=
\begin{cases}
O(\lambda^2\Gamma^{-3}), & y\ll1,\\
O(\lambda^2\Gamma^{-2}), & y\gg1.
\end{cases}
\]

\subsection{15.2 Control knobs}\label{control-knobs}

Candidate controls for lifting a no-go condition or enhancing a go
channel include

\begin{itemize}
\tightlist
\item
  transverse magnetic field
\item
  spin-orbit mixing
\item
  polarization rotation
\item
  strain-induced eigenvector mixing
\item
  cavity-selected submanifold
\item
  spectral selection
\item
  a Raman pathway through another manifold
\end{itemize}

.

\subsection{15.3 General treatment of
strain}\label{general-treatment-of-strain}

Spin-scalar strain does not by itself generate a spin-changing first
moment. It can nevertheless change

\begin{itemize}
\tightlist
\item
  detuning
\item
  orbital eigenvectors
\item
  spin-orbit admixture
\item
  branch linewidth
\item
  selected manifold
\end{itemize}

and therefore enhance or suppress an existing channel. No universal sign
should be assigned to the effect of strain.

\begin{center}\rule{0.5\linewidth}{0.5pt}\end{center}

\section{16. Status of Representative Channels in the Present
Theory}\label{status-of-representative-channels-in-the-present-theory}

\subsection{\texorpdfstring{16.1 NV\(^{-}\) zero-field ZPL
spin-\(\Lambda\)}{16.1 NV\^{}\{-\} zero-field ZPL spin-\textbackslash Lambda}}\label{nv--zero-field-zpl-spin-lambda}

In a simplified parity- or sector-preserving model, a spin-scalar dipole
and a common spin basis can produce an exact sector zero.

In the physical full \(^{3}E\) model, transverse spin--orbit coupling,
nonsecular spin--spin coupling, hyperfine coupling, and related terms
can break the sector. The defensible statement is therefore the
asymptotic dissipative selection rule

\[
M_0=0,
\qquad
K=O(\Gamma^{-2})
\]

. An all-orders exact no-go requires either a reducing symmetry of the
full generator or the vanishing of every moment.

\subsection{16.2 NV transverse-field
channel}\label{nv-transverse-field-channel}

If a weak transverse field is parameterized by
\(\lambda\propto B_\perp\), the simplified model gives

\[
K_{\rm even}=O(B_\perp\Gamma^{-2}),
\qquad
|K_{\rm even}|^2=O(B_\perp^2\Gamma^{-4})
\]

. Even when an odd channel vanishes in a parity-preserving model, the
residual generated by realistic perturbations must be evaluated
separately.

The numerical result currently obtained in the simplified model is

\begin{itemize}
\tightlist
\item
  \(M_0=0\)
\item
  even channel: \(M_1\simeq139\,\mathrm{MHz}\)
\item
  odd channel: \(M_1\simeq0\)
\item
  \(|K|\propto B_\perp\)
\item
  \(K\propto\Gamma^{-2}\)
\end{itemize}

. This is a model-specific result that has not yet been verified with
the full microscopic NV Hamiltonian and is therefore not part of the
universal theorem.

\subsection{\texorpdfstring{16.3 Same-spin
orbital-\(\Lambda\)}{16.3 Same-spin orbital-\textbackslash Lambda}}\label{same-spin-orbital-lambda}

Within the same internal-spin sector, if

\[
d_1^\dagger D^{-1}d_2\neq0
\]

then

\[
K_{12}=O(\Gamma^{-1})
\]

and the channel is a leading-order kernel go. Candidate platforms such
as group-IV centers, orbital doublets, and rare-earth crystal-field
states require their own dipole tensors, phonon relaxation, and
ground-coherence parameters before a practical-go classification can be
made.

\subsection{16.4 Atomic、rare-earth、molecular、quantum-dot
systems}\label{atomicrare-earthmolecularquantum-dot-systems}

The same framework can be used, although the dominant physics differs.

\begin{itemize}
\tightlist
\item
  atomic vapor: Doppler average、collisions、Zeeman manifold
\item
  rare-earth ion: crystal-field/hyperfine states、inhomogeneous
  broadening、long optical lifetime
\item
  molecule: vibronic continuum、Franck--Condon phases、non-Markovian
  bath
\item
  quantum dot: fine-structure splitting、charge noise、cavity and
  polarization selection
\end{itemize}

These systems are classified by configuration tuple rather than by
material name.

\begin{center}\rule{0.5\linewidth}{0.5pt}\end{center}

\section{17. Algorithmic decision tree}\label{algorithmic-decision-tree}

The following is the standard classification algorithm.

\begin{Shaded}
\begin{Highlighting}[]
\NormalTok{INPUT:}
\NormalTok{  H\_g, H\_e, dipole vectors C, jump operators \{L\_mu\},}
\NormalTok{  field/strain/polarization, temperature{-}dependent rates,}
\NormalTok{  probe/control parameters, observable definition}

\NormalTok{1. Build the full zeroth{-}order Liouvillian L\_c and solve rho\^{}(0).}
\NormalTok{2. Restrict the first{-}order problem to the compatible trace{-}zero response space.}
\NormalTok{3. Define the targeted long{-}lived coherence sector S and complementary response sector X.}
\NormalTok{4. Test whether fixed{-}ground optical{-}coherence blocks are invariant.}
\NormalTok{   {-} YES: construct reduced A\_j(z).}
\NormalTok{   {-} NO: retain the full Liouville{-}space block.}
\NormalTok{5. Compute optical dark{-}state dimension: N\_g {-} rank(Omega).}
\NormalTok{6. Compute S\_i and K\_ij where the reduced kernel is valid.}
\NormalTok{7. Construct chi\_full, pathway{-}cut chi\_cut, and delta\_chi\_S.}
\NormalTok{8. Test exact targeted zero:}
\NormalTok{   a. reducing symmetry of the targeted cross{-}transfer;}
\NormalTok{   b. finite Krylov moments;}
\NormalTok{   c. adjugate numerator identity;}
\NormalTok{   d. Laurent/Drazin test at singular points.}
\NormalTok{9. If not exact zero, decompose A = Gamma D + A\_0.}
\NormalTok{10. Determine whether D is full rank.}
\NormalTok{    {-} full rank: compute moment hierarchy and suppression order.}
\NormalTok{    {-} singular: compute protected{-}subspace Schur complement.}
\NormalTok{11. For multiple rates, follow the physical trajectory Gamma\_r(T).}
\NormalTok{12. Compute observable absorption/transmission/fluorescence.}
\NormalTok{13. Distinguish EIT from ATS and population/saturation effects.}
\NormalTok{14. Ensemble{-}average static disorder at the response level.}
\NormalTok{15. Compare contrast/SNR with the experimental threshold.}

\NormalTok{OUTPUT:}
\NormalTok{  optical dark{-}state rank,}
\NormalTok{  exact/asymptotic targeted{-}kernel class and order,}
\NormalTok{  K\_12, K\_21, S\_1, S\_2,}
\NormalTok{  chi\_full, chi\_cut, delta\_chi\_S,}
\NormalTok{  observable EIT contrast,}
\NormalTok{  EIT/ATS label,}
\NormalTok{  parameter margin and uncertainty,}
\NormalTok{  exact assumptions and projectors used.}
\end{Highlighting}
\end{Shaded}

A zero of the total response is not used as a no-go criterion. A no-go
is identified as an identically vanishing targeted coherence feedback.

\section{18. Final theorem package --- corrected finite-dimensional
version}\label{final-theorem-package-corrected-finite-dimensional-version}

\subsection{Theorem 1A --- Optical dark-subspace
rank}\label{theorem-1a-optical-dark-subspace-rank-1}

\[
\mathcal D_{\rm opt}=\ker\Omega,
\qquad
\dim\mathcal D_{\rm opt}
=N_g-\operatorname{rank}\Omega.
\]

This theorem concerns instantaneous optical decoupling and does not by
itself guarantee stationary EIT.

\subsection{Theorem 1B --- Stationary pure Lindblad dark
state}\label{theorem-1b-stationary-pure-lindblad-dark-state}

A pure state \(|D\rangle\langle D|\) is stationary if and only if it is
a common eigenvector of every jump operator and an eigenvector of the
shifted Hamiltonian. An EIT dark state additionally requires
\(\Omega|D\rangle=0\).

\subsection{Theorem 2A --- Full-to-reduced weak-probe
susceptibility}\label{theorem-2a-full-to-reduced-weak-probe-susceptibility}

The full linear response is

\[
\rho^{(1)}
=-\mathcal L_c^\#\mathcal V_p\rho^{(0)}.
\]

. For a closed two-lower-state block, the Schur complement gives

\[
\Xi=S_1-
\frac{\beta K_{12}K_{21}}
{\gamma_g+\beta S_2}.
\]

\subsection{Theorem 2B --- Targeted full-Liouvillian sector-resolved EIT
no-go}\label{theorem-2b-targeted-full-liouvillian-sector-resolved-eit-no-go-1}

Write the first-order response block as

\[
\mathbb M=
\begin{pmatrix}A&B\\C&G_g\end{pmatrix},
\qquad
S_g=G_g-CA^{-1}B
\]

. At a regular point,

\[
\delta\chi_{\mathbb S}
=\chi_{\rm full}-\chi_{\rm cut}^{(\mathbb S)}
=\chi_0d^\dagger A^{-1}BS_g^{-1}CA^{-1}b_p.
\]

If \(\delta\chi_{\mathbb S}\equiv0\) on the specified domain, the net
source-to-detector contribution of the sector \(\mathbb S\) vanishes and
the configuration is an exact \(\mathbb S\)-sector EIT no-go. The
condition \(\chi_{\rm full}=0\) is not a no-go condition and may
represent perfect EIT.

\subsection{Theorem 3 --- Dissipative moment
hierarchy}\label{theorem-3-dissipative-moment-hierarchy}

If \(A_\Gamma=\Gamma D+A_0\), \(D\) is invertible, and the large-rate
expansion converges, then a first nonzero moment \(M_m\) gives

\[
K_\Gamma
=M_m\Gamma^{-m-1}
+O(\Gamma^{-m-2}).
\]

\subsection{Theorem 4 --- Exact finite-dimensional transfer
zero}\label{theorem-4-exact-finite-dimensional-transfer-zero}

For response dimension \(N\),

\[
M_0=\cdots=M_{N-1}=0
\]

is equivalent to finite Krylov orthogonality and to

\[
K_\Gamma=0
\quad
(\Gamma\notin\operatorname{spec}X)
\]

. This conclusion is not restricted to the Neumann-series convergence
domain. For an EIT no-go application, choose \(K_\Gamma\) to be the
targeted sector cross-transfer.

\subsection{Theorem 5 --- Full-generator reducing-symmetry transfer
zero}\label{theorem-5-full-generator-reducing-symmetry-transfer-zero}

If the full Liouvillian or reduced generator has orthogonal reducing
sectors and the source and detector lie in different sectors, the
group-inverse/resolvent transfer vanishes exactly. To obtain an EIT
no-go, define the source and detector for the targeted
coherence-feedback transfer.

\subsection{Theorem 6 --- Singular-damping protected
channel}\label{theorem-6-singular-damping-protected-channel}

Under \(D=D^\dagger\ge0\), the Schur complement on \(\ker D\) determines
the leading response.

\[
K_\Gamma
=p^\dagger P(PA_0P)^{-1}Pc+O(\Gamma^{-1}).
\]

If the leading scalar is nonzero, an \(O(1)\) protected channel remains.

\subsection{Theorem 7 --- Adjugate, analytic identity, and singular
transfer
criterion}\label{theorem-7-adjugate-analytic-identity-and-singular-transfer-criterion}

On the regular set of a connected holomorphic or real-analytic domain,

\[
p^\dagger\operatorname{adj}A(\theta)c\equiv0
\]

gives an exact transfer zero. At a singular point, distinguish Laurent
pole coefficients, finite parts, removable continuations, and the Drazin
inverse; do not interpret \(0/0\) as zero.

\subsection{Corollary --- Symmetric two-branch hopping-rate
map}\label{corollary-symmetric-two-branch-hopping-rate-map}

\[
\dot{(p_X-p_Y)}=-\Gamma_{XY}(p_X-p_Y)
\]

For the symmetric hopping model defined by

\[
\boxed{\gamma_{\rm oc}=\Gamma_{XY}/4}.
\]

\subsection{Corollary --- Weak sector-breaking
control}\label{corollary-weak-sector-breaking-control}

the additional damping of each optical coherence is

\[
K_{12},K_{21}=O(\lambda\Gamma^{-2})
\]

If reciprocal kernels open at the same order through one sector-changing
insertion,

\[
O(\lambda^2\Gamma^{-3})
\quad\text{or}\quad
O(\lambda^2\Gamma^{-2})
\]

. In a nonreciprocal model, evaluate \(K_{12}\) and \(K_{21}\)
separately.

\section{19. Falsifiable predictions and standard
outputs}\label{falsifiable-predictions-and-standard-outputs}

To make the theory experimentally falsifiable, output the following
quantities for every material.

\subsection{19.1 Thermal/damping
collapse}\label{thermaldamping-collapse}

\[
K_{12}(T),\quad K_{21}(T),\quad
\delta\chi_{\mathbb S}(T),\quad C_{\rm EIT}(T)
\]

as functions of the material-specific fast-rate ratio

\[
x(T)=\Gamma_{\rm fast}(T)/\Delta_{\rm sel}
\]

. Use a single universal crossover function only within the model class
for which it has been derived.

\subsection{19.2 External-field opening}\label{external-field-opening}

At small field, measure scaling laws such as

\[
|K|^2\propto B_\perp^2
\]

and the exponent associated with the relevant visibility branch.

\subsection{19.3 Polarization null test}\label{polarization-null-test}

Compare a polarization for which symmetry enforces a zero with one for
which the response is nonzero.

\subsection{19.4 Fast exchange vs static
broadening}\label{fast-exchange-vs-static-broadening}

Test whether two conditions with the same observed linewidth exhibit
different EIT survival.

\subsection{19.5 Material comparison
table}\label{material-comparison-table}

For each configuration, construct

\begin{longtable}[]{@{}
  >{\raggedright\arraybackslash}p{(\columnwidth - 10\tabcolsep) * \real{0.1429}}
  >{\raggedleft\arraybackslash}p{(\columnwidth - 10\tabcolsep) * \real{0.1905}}
  >{\raggedleft\arraybackslash}p{(\columnwidth - 10\tabcolsep) * \real{0.1905}}
  >{\raggedright\arraybackslash}p{(\columnwidth - 10\tabcolsep) * \real{0.1429}}
  >{\raggedleft\arraybackslash}p{(\columnwidth - 10\tabcolsep) * \real{0.1905}}
  >{\raggedright\arraybackslash}p{(\columnwidth - 10\tabcolsep) * \real{0.1429}}@{}}
\toprule\noalign{}
\begin{minipage}[b]{\linewidth}\raggedright
configuration
\end{minipage} & \begin{minipage}[b]{\linewidth}\raggedleft
dark rank
\end{minipage} & \begin{minipage}[b]{\linewidth}\raggedleft
first nonzero moment
\end{minipage} & \begin{minipage}[b]{\linewidth}\raggedright
targeted sector class
\end{minipage} & \begin{minipage}[b]{\linewidth}\raggedleft
predicted contrast
\end{minipage} & \begin{minipage}[b]{\linewidth}\raggedright
evidence status
\end{minipage} \\
\midrule\noalign{}
\endhead
\bottomrule\noalign{}
\endlastfoot
material/channel A & & & exact/asymptotic/go & & measured/predicted \\
\end{longtable}

.

\begin{center}\rule{0.5\linewidth}{0.5pt}\end{center}

\section{20. Current Theoretical
Limitations}\label{current-theoretical-limitations}

The following lie outside the integrated theory or require additional
proof.

\begin{enumerate}
\def\labelenumi{\arabic{enumi}.}
\tightlist
\item
  A finite-moment exact theorem for an arbitrary infinite-dimensional
  non-Markovian bath
\item
  A simple kernel formula under arbitrary strong-probe saturation
\item
  Collective EIT for which propagation, cavity feedback, or many-body
  interactions cannot be neglected
\item
  Inferring an all-orders exact zero from a numerical zero
\item
  Inferring observable EIT directly from \(K\neq0\)
\item
  Classifying EIT, CPT, STIRAP, Raman transfer, and quantum memory under
  one identical condition
\item
  Assigning one go/no-go label to an entire material family
\item
  Defining a unique full-Liouvillian EIT no-go without specifying the
  coherence-sector partition
\item
  Interpreting a zero of the total susceptibility as an EIT no-go
\end{enumerate}

STIRAP requires not only a nonzero Raman kernel but also a
time-dependent dark eigenstate, a counterintuitive pulse sequence, and
an adiabatic gap.

\begin{center}\rule{0.5\linewidth}{0.5pt}\end{center}

\section{21. Theory-completion audit and remaining validation
tasks}\label{theory-completion-audit-and-remaining-validation-tasks}

\subsection{21.1 Completion status}\label{completion-status}

The term ``completed'' in this version refers only to the stated
assumptions of finite dimensionality, Markovian GKSL dynamics, a
stationary rotating frame, weak probing, and an analytic parameter
domain.

\begin{longtable}[]{@{}
  >{\raggedright\arraybackslash}p{(\columnwidth - 4\tabcolsep) * \real{0.3333}}
  >{\raggedright\arraybackslash}p{(\columnwidth - 4\tabcolsep) * \real{0.3333}}
  >{\raggedright\arraybackslash}p{(\columnwidth - 4\tabcolsep) * \real{0.3333}}@{}}
\toprule\noalign{}
\begin{minipage}[b]{\linewidth}\raggedright
requested item
\end{minipage} & \begin{minipage}[b]{\linewidth}\raggedright
status in v6.2
\end{minipage} & \begin{minipage}[b]{\linewidth}\raggedright
exact scope
\end{minipage} \\
\midrule\noalign{}
\endhead
\bottomrule\noalign{}
\endlastfoot
Dark-state theorem & \textbf{completed} & Separates optical nullity from
a stationary Lindblad dark state and gives necessary and sufficient
conditions \\
Exact finite-dimensional transfer-zero proof & \textbf{completed} &
Extended to every invertible \(\Gamma\) using Cayley--Hamilton, Krylov,
and adjugate arguments \\
Full-Liouvillian EIT no-go definition & \textbf{completed} & Defined as
the sector-resolved condition \(\delta\chi_{\mathbb S}\equiv0\), not as
a zero of the total susceptibility \\
Adjugate singularities & \textbf{completed} & Separates the regular
domain, Laurent poles, finite parts, removable continuations, and Drazin
treatment \\
Weak-probe susceptibility & \textbf{completed} & Derived from the
full-Liouvillian group inverse and proved in scalar/matrix
Schur-complement form \\
Singular ground block & \textbf{completed} & Uses \(S_g=G_g-CA^{-1}B\)
without requiring \(G_g^{-1}\) \\
\(\Gamma_{XY}\) mapping & \textbf{completed for the stated convention} &
For a symmetric population-imbalance rate,
\(\gamma_{\rm oc}=\Gamma_{XY}/4\) \\
Scope of \(D>0\) & \textbf{completed} & Separates invertibility, strict
accretivity, and the Hermitian singular theorem \\
Full-generator symmetry & \textbf{completed as a transfer theorem} &
Audits every jump operator, the steady-state projector, source, and
detector, not only the Hamiltonian \\
Analytic identity condition & \textbf{completed} & Restricted to
holomorphic or real-analytic domains, without overextending to smooth
dependence \\
Sector dependence & \textbf{completed} & States the \(\mathbb S\)
dependence, block-preserving basis invariance, and dependence under
sector mixing \\
Textbook \(\Lambda\) recovery & \textbf{completed} & Recovers the
standard three-level susceptibility exactly \\
Mishina-type recovery & \textbf{completed} & Recovers the diagonal
multi-excited resolvent sum and inhomogeneous averaging \\
Full/reduced unit test & \textbf{completed} & Compares a physical
three-level GKSL Liouvillian with the Schur formula to machine
precision \\
\end{longtable}

\subsection{21.2 Theoretical conditions for proceeding to numerical
calculations}\label{theoretical-conditions-for-proceeding-to-numerical-calculations}

The present version satisfies the following conditions.

\begin{enumerate}
\def\labelenumi{\arabic{enumi}.}
\tightlist
\item
  Optical dark states and stationary dark states are not conflated.
\item
  An exact transfer zero does not depend on the convergence domain of
  the large-rate series.
\item
  An ordinary inverse is not used for a singular Liouvillian.
\item
  The reduced susceptibility is identified as the Schur complement of
  the full first-order response.
\item
  A Schur complement valid for singular \(G_g\) is used.
\item
  The coefficient relating a population-hopping rate to an
  optical-coherence damping rate is derived from the jump convention.
\item
  The use of positivity assumptions is restricted to the locations where
  they are required.
\item
  The symmetry theorem is formulated as a reducing-sector transfer
  theorem for the full generator.
\item
  A zero of the total response is separated from an EIT no-go.
\item
  The targeted sector \(\mathbb S\), \(\chi_{\rm cut}^{(\mathbb S)}\),
  and \(\delta\chi_{\mathbb S}\) are explicitly specified.
\item
  The sector dependence and block-preserving basis invariance of
  \(\delta\chi_{\mathbb S}\) are stated explicitly.
\item
  The textbook \(\Lambda\) model and a Mishina-type diagonal multilevel
  response are recovered as special cases.
\item
  Unit tests comparing the full GKSL Liouvillian and reduced formula
  have been completed to machine precision.
\end{enumerate}

Accordingly, there is no need to prioritize further abstract
generalization before proceeding to material-specific numerical
calculations.

\subsection{21.3 Completed full-Liouvillian and reduced-formula unit
tests}\label{completed-full-liouvillian-and-reduced-formula-unit-tests}

The core regression script \texttt{eit\_unit\_tests\_v6\_2.py}
independently verifies the following results using double-precision
complex arithmetic.

\subsubsection{\texorpdfstring{Test U1 --- Physical three-level GKSL
Liouvillian vs reduced \(\Lambda\)
formula}{Test U1 --- Physical three-level GKSL Liouvillian vs reduced \textbackslash Lambda formula}}\label{test-u1-physical-three-level-gksl-liouvillian-vs-reduced-lambda-formula}

Using the basis \((|g_1\rangle,|g_2\rangle,|e\rangle)\),

\[
\gamma_{31}=0.73,
\quad
\gamma_{21}=0.041,
\quad
\Delta_p=0.37,
\quad
\delta_2=-0.083,
\quad
\Omega_c=1.16
\]

was adopted. A full GKSL Liouvillian containing excited-state decay
jumps and a \(|g_2\rangle\to|g_1\rangle\) relaxation jump was
constructed, and the zeroth-order steady state and trace-zero
first-order response were solved directly.

\[
\rho_{eg_1,\rm full}^{(1)}
=-0.09734305655387952
+0.06276778568576767i,
\]

\[
\rho_{eg_1,\rm reduced}^{(1)}
=-0.09734305655387981
+0.06276778568576774i.
\]

\[
\boxed{
\varepsilon_{\rm rel}=2.996\times10^{-16}
}
\]

Thus, the full-Liouvillian weak-probe solution and the reduced formula
of Section 9.6 agree to machine precision.

\subsubsection{Test U2 --- Textbook perfect-EIT
limit}\label{test-u2-textbook-perfect-eit-limit}

\[
\gamma_{21}=\delta_2=0,
\qquad
\beta>0
\]

one obtains

\[
\Xi_{\rm full}=0,
\qquad
\Xi_{\rm cut}=A^{-1},
\qquad
\delta\Xi_{\mathbb S}=-A^{-1}\neq0.
\]

For the numerical example \(A=0.9-0.2i\),

\[
\Xi_{\rm cut}=1.0588235294+0.2352941176i.
\]

This fixes the convention so that a zero of the total susceptibility is
not misidentified as a no-go.

\subsubsection{Test U3 --- Mishina-type diagonal multi-excited-state
recovery}\label{test-u3-mishina-type-diagonal-multi-excited-state-recovery}

For a random complex diagonal generator with four excited branches, the
kernel sum of Section 9.7 was compared with direct inversion of the
\((N_e+1)\)-dimensional block.

\[
\Xi_{\rm kernel}
=1.9240573122129712-0.06390491266120876i,
\]

\[
\Xi_{\rm direct}
=1.9240573122129716-0.06390491266120868i,
\]

\[
\boxed{
\varepsilon_{\rm rel}=2.347\times10^{-16}
}.
\]

\subsubsection{Test U4 --- Block-preserving basis
invariance}\label{test-u4-block-preserving-basis-invariance}

For a random \(3\oplus2\) block system, independent unitaries
\(U_X\oplus U_S\) were applied.

\[
\boxed{
\frac{|\delta\chi_{\mathbb S}'-\delta\chi_{\mathbb S}|}
{\max(1,|\delta\chi_{\mathbb S}|)}
=2.734\times10^{-16}
}.
\]

\subsubsection{Test U5 --- Sector-mixing
dependence}\label{test-u5-sector-mixing-dependence}

After a unitary transformation mixing \(X\) and \(S\), the coordinate
cut was redefined, yielding

\[
\delta\chi_{\rm before}
=-0.0049367+0.1224286i,
\]

\[
\delta\chi_{\rm after\ recut}
=-0.2348257-0.0281825i,
\]

\[
|\Delta|=0.2748321.
\]

This numerically confirms that \(\delta\chi_{\mathbb S}\) is invariant
under basis changes within the sectors but depends on the physical
definition of the sector.

\subsubsection{Existing algebraic regression
tests}\label{existing-algebraic-regression-tests}

\begin{longtable}[]{@{}
  >{\raggedright\arraybackslash}p{(\columnwidth - 2\tabcolsep) * \real{0.4286}}
  >{\raggedleft\arraybackslash}p{(\columnwidth - 2\tabcolsep) * \real{0.5714}}@{}}
\toprule\noalign{}
\begin{minipage}[b]{\linewidth}\raggedright
test
\end{minipage} & \begin{minipage}[b]{\linewidth}\raggedleft
result
\end{minipage} \\
\midrule\noalign{}
\endhead
\bottomrule\noalign{}
\endlastfoot
exact Krylov transfer zero & residual \(0\) \\
one-insertion hierarchy & fitted slope \(-2.0000000000000004\) \\
symmetric hopping map & \(\gamma_{\rm oc}-\Gamma_{XY}/4=0\) \\
scalar weak-probe Schur complement & relative residual
\(1.62\times10^{-16}\) \\
matrix Schur complement with singular \(G_g\) & residual \(0\) \\
Drazin finite part & regularization error \(2.50\times10^{-4}\) at
finite shift \\
symmetry-projector transfer zero & residual \(0\) \\
\end{longtable}

These tests are consistency checks of the finite-dimensional theorems
and solver architecture; they do not replace validation of the
material-specific Hamiltonian, jump set, rate convention, or observable
model.

\subsection{21.4 Remaining material-specific numerical and model
validation}\label{remaining-material-specific-numerical-and-model-validation}

\subsubsection{Priority A --- Material-solver regression and observable
validation}\label{priority-a-material-solver-regression-and-observable-validation}

The core full/reduced unit tests are complete. When moving to a material
model, maintain the following as regression tests for every parameter
set.

\begin{itemize}
\tightlist
\item
  Domain of agreement between the full-Liouvillian group/Drazin solve
  and the reduced block
\item
  Reconstruction error among \(\chi_{\rm full}\),
  \(\chi_{\rm cut}^{(\mathbb S)}\), and \(\delta\chi_{\mathbb S}\)
\item
  Krylov exact-zero and Laurent/singular-point tests
\item
  Nonnormal resolvent norm and condition number
\item
  EIT/ATS pole-residue classifier
\item
  Audit of truncation dimension, steady-state uniqueness, and trace
  preservation
\end{itemize}

\subsubsection{Priority B --- NV microscopic
audit}\label{priority-b-nv-microscopic-audit}

\begin{itemize}
\tightlist
\item
  full six-level \(^{3}E\) Hamiltonian
\item
  physical complex dipole vectors
\item
  Mapping of \(k_{X\leftrightarrow Y}(T)\) from the jump convention
\item
  full-generator symmetry audit table
\item
  Verification of \(M_0,M_1\) and the predicted suppression order
\item
  transverse-field／strain symmetry breaking
\item
  Comparison of
  \(\chi_{\rm full},\chi_{\rm cut}^{(\mathbb S)},\delta\chi_{\mathbb S}\)
\end{itemize}

\subsubsection{Priority C --- Cross-material
validation}\label{priority-c-cross-material-validation}

Using the same solver and observable normalization, demonstrate

\begin{enumerate}
\def\labelenumi{\arabic{enumi}.}
\tightlist
\item
  exact targeted no-go case
\item
  asymptotic no-go case
\item
  a leading-order or protected go case
\item
  experimentally established EIT benchmark
\item
  field/polarization tunable escape case
\end{enumerate}

.

\textbf{Assessment:} Within the stated assumptions, the principal
mathematical inconsistencies in the general theory have been corrected.
The main remaining work is not the addition of further abstract
theorems, but the validation of microscopic inputs and solvers and the
production of material-specific numerical predictions.

\section{22. Recommended Central Claim of the
Paper}\label{recommended-central-claim-of-the-paper}

\subsection{English}\label{english}

\begin{quote}
For any finite-dimensional Markovian open quantum emitter admitting a
stationary rotating-frame weak-probe response, and for a physically
specified long-lived sector S, the S-mediated probe-response
modification is a finite-dimensional transfer function. Exact S-sector
no-go corresponds to an identically vanishing sector-resolved
correction, not to a vanishing total susceptibility. The framework
recovers the textbook three-level Lambda susceptibility and the diagonal
multi-excited-state Mishina-type response. Fast full-rank damping
produces a hierarchy whose suppression order is fixed by the first
nonvanishing resolvent moment, while observable EIT requires the full
Liouvillian visibility budget.
\end{quote}

\subsection{Japanese version}\label{japanese-version}

\begin{quote}
For a finite-dimensional Markovian open quantum emitter admitting a
weak-probe response in a stationary rotating frame, the probe-response
modification mediated by a physically specified long-lived sector \(S\)
is a finite-dimensional transfer function. If this sector-resolved
correction vanishes identically, the system is an exact \(S\)-sector
no-go; a zero of the total susceptibility is not a no-go condition. The
framework recovers the textbook three-level \(\Lambda\) susceptibility
and a Mishina-type diagonal multi-excited-state response as special
cases. Under strong full-rank damping, the first nonzero resolvent
moment determines the suppression order, while observable EIT requires a
separate visibility condition based on the full Liouvillian.
\end{quote}

\begin{center}\rule{0.5\linewidth}{0.5pt}\end{center}

\section{23. Final Conclusion}\label{final-conclusion}

The theory has been extended from an NV-specific no-go argument to a
material-independent framework containing the following elements.

\begin{enumerate}
\def\labelenumi{\arabic{enumi}.}
\tightlist
\item
  Hilbert-space dark-state rank
\item
  basis-invariant matrix resolvent kernel
\item
  full Liouville-space pathway-resolved response correction
\item
  a moment hierarchy for full-rank damping
\item
  exact Krylov/adjugate/symmetry transfer-zero criteria
\item
  protected channels under singular damping
\item
  multiple-rate and external-field phase diagram
\item
  observable-specific EIT visibility
\item
  separation of static broadening from fast dynamics
\item
  a channel-resolved taxonomy
\end{enumerate}

The theory can therefore be applied to arbitrary materials. The final
classification, however, is performed not by material name but for each
configuration after fixing the material, lower-state pair, excited
manifold, polarizations, external fields, temperature, observable, and
targeted coherence sector. An exact EIT no-go is
\(\delta\chi_{\mathbb S}\equiv0\), not \(\chi_{\rm full}=0\). This is
the strongest generalization that does not overstate the theory.

\begin{center}\rule{0.5\linewidth}{0.5pt}\end{center}

\section{Appendix A. Notation}\label{appendix-a.-notation}

\begin{longtable}[]{@{}
  >{\raggedright\arraybackslash}p{(\columnwidth - 2\tabcolsep) * \real{0.5000}}
  >{\raggedright\arraybackslash}p{(\columnwidth - 2\tabcolsep) * \real{0.5000}}@{}}
\toprule\noalign{}
\begin{minipage}[b]{\linewidth}\raggedright
symbol
\end{minipage} & \begin{minipage}[b]{\linewidth}\raggedright
meaning
\end{minipage} \\
\midrule\noalign{}
\endhead
\bottomrule\noalign{}
\endlastfoot
\(\mathcal H_g,\mathcal H_e\) & lower and excited manifolds \\
\(\Omega\) & instantaneous optical coupling map \\
\(C\) & matrix of dipole coupling vectors \\
\(A(z)\) & optical-coherence generator \\
\(G=A^{-1}\) & resolvent/Green operator \\
\(K_{ab}\) & off-diagonal coherent-response kernel \\
\(S_a\) & diagonal optical response \\
\(D\) & fast damping-shape operator \\
\(\Gamma\) & fast optical-coherence damping scale \\
\(M_n=p^\dagger X^n\nu\) & resolvent moment \\
\(Q\) & reducing symmetry operator \\
\(P=\operatorname{Proj}(\ker D)\) & protected-subspace projector \\
\(\gamma_g\) & complex lower-coherence decay/detuning \\
\(\beta=|\Omega_c|^2/4\) & control intensity parameter \\
\(\Xi\) & normalized full local probe response \\
\(\chi_{\rm cut}^{(\mathbb S)}\) & algebraic counterfactual response
obtained by cutting the cross-feedback with the specified sector \\
\(\delta\chi_{\mathbb S}\) &
\(\chi_{\rm full}-\chi_{\rm cut}^{(\mathbb S)}\), sector-mediated
response modification \\
\(C_{\rm abs}\) & specified relative absorption contrast \\
\end{longtable}

\section{Appendix B. Proof-repair
ledger}\label{appendix-b.-proof-repair-ledger}

The principal points corrected and fixed in this version are as follows.

\begin{itemize}
\tightlist
\item
  The dark-state rank is restricted to an optical-nullity theorem, while
  the full Lindblad necessary and sufficient conditions for a stationary
  pure state are retained.
\item
  The moment symbol is standardized as \(\nu=D^{-1}c\).
\item
  An exact transfer zero is proved through the Krylov/Cayley--Hamilton
  equivalence.
\item
  The adjugate criterion is restricted to an analytic regular domain.
\item
  At singular points, finite parts, removable continuations, and
  ordinary inverses are distinguished.
\item
  The weak-probe formula is derived from the full-Liouvillian
  group-inverse response.
\item
  The formulation is corrected to use \(S_g=G_g-CA^{-1}B\), without
  requiring the inverse of the scalar ground block.
\item
  The full-Liouvillian EIT no-go is corrected from \(\chi_{\rm full}=0\)
  to \(\delta\chi_{\mathbb S}\equiv0\).
\item
  The ideal \(\Lambda\) system is added as a counterexample, confirming
  that a zero total response may represent perfect EIT.
\item
  The \(\Gamma_{XY}/4\) mapping is verified explicitly for symmetric
  two-branch hopping.
\item
  The locations requiring \(D>0\) are separated from those requiring
  only invertibility or strict accretivity.
\item
  The symmetry theorem is positioned as a generic transfer-zero theorem,
  and its use for an EIT no-go is restricted to the targeted
  cross-transfer.
\item
  Reciprocal and nonreciprocal kernels are distinguished in weak
  sector-breaking scaling.
\end{itemize}

\section{Appendix C. Provenance and integration
status}\label{appendix-c.-provenance-and-integration-status}

This integrated version is based on the v6.0 theory package and
incorporates into the v6.1 proof-corrected formulation the sector
dependence, recovery of standard models, Markovian scope, and
full/reduced unit tests.

The main integrated elements are as follows.

\begin{itemize}
\tightlist
\item
  Excited-confined jumps in the optical-coherence block are treated as
  full-rank damping.
\item
  Exact dark states, coherent pathways, and observable EIT are
  separated.
\item
  Exact transfer zeros are classified by finite moments, Krylov
  orthogonality, adjugate identities, and reducing symmetries.
\item
  An EIT no-go is defined as an identically vanishing pathway-cut
  correction \(\delta\chi_{\mathbb S}\) for a physically specified
  sector \(\mathbb S\).
\item
  A zero of the total susceptibility is excluded from the no-go
  condition.
\item
  A singular ground-coherence block is handled by the \(A\)-side Schur
  complement.
\item
  Multiple rates, singular damping, and static disorder are treated
  within the main scope, while non-Markovian extensions are stated
  explicitly to be outside it.
\item
  No-go/go classifications are made at the configuration/channel level
  rather than the material level.
\end{itemize}

\textbf{Integration status:} Within the stated finite-dimensional
Markovian stationary weak-probe assumptions, the theory has been
corrected into an integrated form suitable for numerical calculations.
Material-specific conclusions must be finalized only after numerical
validation of each Hamiltonian, dipole operator, jump operator, rate,
and observable model.

\end{document}
